% GENERATED by convert_sources.py from docs/spec_power.md. Do not edit by hand.
\subsection{Main PCB: power section circuit specification (v1.2)}

\textbf{Changelog}

\begin{itemize}
\tightlist
\item
  \textbf{v1.2 (2026-11-24): matches the schematic after ECO-001 and ECO-002} (\texttt{docs/\allowbreak{}eco\_\allowbreak{}001.md}, \texttt{docs/\allowbreak{}eco\_\allowbreak{}002.md}, \texttt{docs/\allowbreak{}review/\allowbreak{}release\_\allowbreak{}review.md}).

  \begin{itemize}
  \tightlist
  \item
    Boost compensation: R903 100 k, C915 15 nF.
  \item
    R907/D901 fitted; firmware PVDD soft start.
  \item
    BQ25792 POR/watchdog ICHG is \textbf{2 A} (S3).
  \item
    New DNP C721 and the rule for fitting C720 (S2).
  \item
    JP803 default is now the pack NTC (S6).
  \item
    Balance-lead assembly check (S7).
  \item
    CHG\_\allowbreak{}QON is routed to the UI button through the FFC (S4).
  \item
    Firmware protection contract (S1).
  \end{itemize}
\item
  v1.1: Q901 removed; PVDD 18.0 V; C920 footprint CP\_\allowbreak{}Elec\_\allowbreak{}10x10.
\item
  v1.0: initial release.
\end{itemize}

Status: 2026-11-24. Scope: sheets \texttt{usb\_\allowbreak{}power}, \texttt{charger}, \texttt{battery}, \texttt{rails} of the main board. This document is written so the schematic can be drawn \textbf{without datasheets}: every pin, value, net and footprint is below. Companion spec: \texttt{docs/\allowbreak{}spec\_\allowbreak{}signal.md} (MCU, amps, UI).

\textbf{Conventions} (aligned with \texttt{spec\_\allowbreak{}signal.md})

\begin{itemize}
\tightlist
\item
  Net names in \texttt{CODE} are exact. Nets marked \textbf{(global)} cross sheets; all others are local labels.
\item
  RefDes ranges: \texttt{usb\_\allowbreak{}power} 6xx, \texttt{charger} 7xx, \texttt{battery} 8xx, \texttt{rails} 9xx (boost 90x--92x, buck 95x). The signal sheets own 100--599.
\item
  \textbf{B} = JLC Basic, \textbf{E} = JLC Extended (JLC parts API, 2026-11-24). ``Stock sym/fp'' = present in the KiCad 10.0.4 stock libraries (checked locally).
\item
  Default passive footprints: \texttt{Resistor\_\allowbreak{}SMD:\allowbreak{}R\_\allowbreak{}0603\_\allowbreak{}1608Metric}, \texttt{Capacitor\_\allowbreak{}SMD:\allowbreak{}C\_\allowbreak{}0603\_\allowbreak{}1608Metric}, \texttt{…C\_\allowbreak{}0805\_\allowbreak{}2012Metric}, \texttt{…C\_\allowbreak{}1206\_\allowbreak{}3216Metric}, \texttt{Resistor\_\allowbreak{}SMD:\allowbreak{}R\_\allowbreak{}1206\_\allowbreak{}3216Metric}. Symbols \texttt{Device:\allowbreak{}R}, \texttt{Device:\allowbreak{}C}, \texttt{Device:\allowbreak{}C\_\allowbreak{}Polarized}, \texttt{Device:\allowbreak{}L}, \texttt{Device:\allowbreak{}LED}, \texttt{Device:\allowbreak{}D}, \texttt{Device:\allowbreak{}D\_\allowbreak{}Zener}, \texttt{Device:\allowbreak{}Fuse}, \texttt{Device:\allowbreak{}Thermistor\_\allowbreak{}NTC}.
\item
  ``DNP'' = footprint placed, part not fitted (set the KiCad DNP attribute).
\item
  Datasheets used (local cache: \texttt{/\allowbreak{}private/\allowbreak{}tmp/\allowbreak{}claude-501/\allowbreak{}-Users-frankwalsh-Documents-vibecoding/\allowbreak{}8f7148fa-1ea7-42b7-910f-60a1668ad40c/\allowbreak{}scratchpad/\allowbreak{}datasheets/\allowbreak{}} and \texttt{…/\allowbreak{}datasheets/\allowbreak{}pw/\allowbreak{}}):
  BQ25792 SLUSDG1C (LCSC copy = TI Chinese edition ZHCSMD3C); BQ25798 SLUSDV2C (English, same pinout/sections); BQ2579x EVM guide SLUUCB5E; TPS61178 SLVSDA7E; TPS62933 SLUSEA4D; WCH ``CH224 手册1'' V2.1; Ricoh R5460x2xx EA-165-160603; HYCON HY2213 DS-V05; NCE NCEP3065QU V6.0; Littelfuse 451/453; MDD SMF24A; UMW USBLC6-2SC6; Amass XT30PW-M.
\end{itemize}

\subsubsection{Decisions and changes vs.~the brief (read first)}

\begin{enumerate}
\def\labelenumi{\arabic{enumi}.}
\tightlist
\item
  \textbf{PVDD = 18.0 V nominal} (lead decision: 8 Ω woofers, Dayton ND65-8; 4 Ω tweeters). Boost OVP is 21 V typ (21.5 V max), well under the TAS5805M's 26.4 V maximum.
\item
  \textbf{PVDD\_\allowbreak{}TRIM uses a diode-OR injection} (R907 68 k + D901 1N4148W, \textbf{fitted} since ECO-001), giving \textbf{≈12.2--18.0 V} (SPICE). With the GPIO Hi-Z \textbf{or low}, PVDD stays at 18.0 V. There is no DC path from \texttt{PVDD\_\allowbreak{}TRIM} to GND or to a rail (only the RC capacitors). A high trim voltage lowers PVDD; the trim cannot raise it except by \textless{} 0.1 V of D901 reverse leakage at ≤ 85 °C. Firmware also uses it as a PVDD soft start (§4.1.4).
\item
  \textbf{The TPS61178 load-disconnect FET is NOT used} (DISDRV left open, as the TPS61178 pin table allows). The amp sheets each carry 470 µF + 4×10 µF, and this sheet adds 470 µF, so PVDD holds ≈1.5 mF. TI requires COUT2 ≤ 10 × COUT1 behind the disconnect FET, and the FET must turn on within 3 ms (§9.2.4.2, §9.2.4.4.3). About 40× is not achievable, and start-up would risk hiccup. Consequences, all accepted:

  \begin{itemize}
  \tightlist
  \item
    With BOOST\_\allowbreak{}EN low, PVDD ≈ VSYS − 0.5 V (6--8 V) through the high-side body diode. The amps sit in PDN shutdown, drawing µA.
  \item
    At battery connect, the bulk charges through Q701, the ship FET, acting as a source follower. The inrush is \textbf{not bounded by design}: 2.5--9.7 A estimated (release review S2). C720 (gate--drain) and C721 (gate--source), both DNP, can bound it, subject to the C720 fitting rule in §2.1.
  \item
    A PVDD short is cleared by the BQ IBAT\_\allowbreak{}OCP 9.3 A ship-FET trip and by F801, not by the boost's hiccup mode. \textbf{The trip needs firmware} (SFET\_\allowbreak{}PRESENT = 1, EN\_\allowbreak{}BATOC = 1, watchdog handled; §7). With blank or hung firmware, only the R5460 (≈40 A for 12 ms) and F801 remain (release review S1).
  \end{itemize}
\item
  \textbf{A ship FET is added on the charger's BAT path} (Q701, NCEP3065QU, the same LCSC part as the protection FETs, so no new line item). The BQ2579x \textbf{battery over-current protection (IBAT\_\allowbreak{}OCP 9.3 A) only works with a ship FET fitted} (BQ25798 §7.3.13.1.7). It also gives a true ship/shutdown mode and a QON ``power-on'' button.
\item
  \textbf{HY2213-BB3A is a gate driver, not a bleeder.} Its OUT pin drives an external N-FET plus bleed resistor (HY2213 §9). Each cell gets AO3400A + 2× 120 Ω.
\item
  \textbf{The R5460 is a low-side protector.} Cell negative (\texttt{BATT\_\allowbreak{}N}) is \textbf{not} board GND: the protection FETs sit between \texttt{BATT\_\allowbreak{}N} and \texttt{GND}. Copper and connectors must respect this (§3, §6).
\item
  \textbf{Balance lead: only the middle tap is used by default.} Pins 1 (B−) and 3 (B+) of the JST-XH go through open solder jumpers. If a pack has its own PCM, tying its balance B−/B+ to our main path would let current bypass the pack's protection FETs through the thin balance wires. A 1 MΩ/1 MΩ divider biases the mid node, so the protector and balancers also work with the balance lead unplugged.
\item
  \textbf{CH224Q is isolated from the shared I2C bus when VBUS is absent} (Q601/Q602). Its CFG2/SCL and CFG3/SDA have internal 7--15 kΩ pull-ups to its internal rail (CH224 §8.2.1). Unpowered, that rail sits near 0 V, which would drag the 3.3 V bus to ≈2.2 V, below the ESP32/TAS5805M V\_\allowbreak{}IH.
\item
  \textbf{I2C pull-ups, the CHG\_\allowbreak{}INT pull-up and the BOOST\_\allowbreak{}EN pull-down are not on the power sheets.} They are R104/R105/R106/R107 on the \texttt{mcu} sheet. Do not duplicate them.
\item
  \textbf{Charger switching frequency is 750 kHz with 2.2 µH} (PROG = 8.2 kΩ, 2S). TI: 750 kHz works only with 2.2 µH, 1.5 MHz only with 1 µH (BQ25798 §8.2.2.2).
\end{enumerate}

\begin{center}\rule{0.5\linewidth}{0.5pt}\end{center}

\subsubsection{0. Parts with no stock KiCad symbol and/or footprint}

{\def\LTcaptype{none} % do not increment counter
\begin{longtable}[]{@{}
  >{\raggedright\arraybackslash}p{(\linewidth - 4\tabcolsep) * \real{0.3333}}
  >{\raggedright\arraybackslash}p{(\linewidth - 4\tabcolsep) * \real{0.3333}}
  >{\raggedright\arraybackslash}p{(\linewidth - 4\tabcolsep) * \real{0.3333}}@{}}
\toprule\noalign{}
\begin{minipage}[b]{\linewidth}\raggedright
Part
\end{minipage} & \begin{minipage}[b]{\linewidth}\raggedright
Symbol
\end{minipage} & \begin{minipage}[b]{\linewidth}\raggedright
Footprint
\end{minipage} \\
\midrule\noalign{}
\endhead
\bottomrule\noalign{}
\endlastfoot
BQ25792RQMR & Use stock \textbf{\texttt{Battery\_\allowbreak{}Management:\allowbreak{}BQ25798}}, set Value = \texttt{BQ25792RQMR}. \textbf{The pinouts are identical}: BQ25792 Table 7-1 and BQ25798 Table 5-1 list the same 29 pins with the same numbers and names, both in the RQM VQFN-29. The KiCad BQ25798 symbol's pins (1 STAT \ldots{} 29 PMID) match both tables pin for pin (checked). The only text difference: the BQ25792 SDRV note allows the 1 nF to go to GND \emph{or} BAT when no ship FET is used, which doesn't apply here. & Stock \texttt{Package\_\allowbreak{}DFN\_\allowbreak{}QFN:\allowbreak{}Texas\_\allowbreak{}RQM0029A\_\allowbreak{}VQFN-29\_\allowbreak{}4x4mm\_\allowbreak{}P0.4mm} \\
TPS61178RNWR & \textbf{Custom} \texttt{power:\allowbreak{}TPS61178} (pin table §4.1.1) & \textbf{Custom} \texttt{power:\allowbreak{}Texas\_\allowbreak{}RNW0013A\_\allowbreak{}VQFN-HR-13\_\allowbreak{}3x3.5mm\_\allowbreak{}P0.5mm} (geometry §4.1.2) \\
CH224Q & \textbf{Custom} \texttt{power:\allowbreak{}CH224Q}. Easiest: copy stock \texttt{Interface\_\allowbreak{}USB:\allowbreak{}CH224K} (same pin numbers 1--11) and rename pin 1 → \texttt{VHV}, pin 2 → \texttt{CFG2/\allowbreak{}SCL}, pin 3 → \texttt{CFG3/\allowbreak{}SDA} (pin table §1.2) & Stock \texttt{Package\_\allowbreak{}DFN\_\allowbreak{}QFN:\allowbreak{}Texas\_\allowbreak{}DSQ0010A\_\allowbreak{}WSON-10-1EP\_\allowbreak{}2x2mm\_\allowbreak{}P0.4mm\_\allowbreak{}EP0.9x1.5mm} (pad 11 = EP = GND). CH224Q DFN10: body 2.0×2.0, pitch 0.4, lead 0.2×0.3, EP 1.5×0.9 (CH224 §9.2). The TI DSQ pattern has the same geometry rotated 90°, with counter-clockwise numbering from pin 1 like the CH224Q \\
R5460N212AF & \textbf{Custom} \texttt{power:\allowbreak{}R5460x2} (pin table §3.2) & Stock \texttt{Package\_\allowbreak{}TO\_\allowbreak{}SOT\_\allowbreak{}SMD:\allowbreak{}SOT-23-6} \\
HY2213-BB3A & \textbf{Custom} \texttt{power:\allowbreak{}HY2213} (pin table §3.3) & Stock \texttt{Package\_\allowbreak{}TO\_\allowbreak{}SOT\_\allowbreak{}SMD:\allowbreak{}SOT-23-6} \\
NCEP3065QU & Stock \texttt{Transistor\_\allowbreak{}FET:\allowbreak{}CSD17581Q3A} (= \texttt{Q\_\allowbreak{}NMOS\_\allowbreak{}SSSGD\_\allowbreak{}AvalancheRated}: pins 1-2-3 S, 4 G, 5 D), Value = \texttt{NCEP3065QU} & Stock \texttt{Package\_\allowbreak{}SON:\allowbreak{}VSON-8\_\allowbreak{}3.3x3.3mm\_\allowbreak{}P0.65mm\_\allowbreak{}NexFET} (pad 5 = drain tab). NCE DFN3.3×3.3-8L: pitch 0.65, lead 0.30, drain pad 2.49×1.88 (NCE p.5), which matches this pattern \\
0451010.MRL (2410 fuse) & Stock \texttt{Device:\allowbreak{}Fuse} & \textbf{Custom} \texttt{power:\allowbreak{}Fuse\_\allowbreak{}2410\_\allowbreak{}6125Metric}: 2 SMD rect pads 1.96 mm (x) × 3.15 mm (y), centres at x = ±2.455 mm (gap 2.95 mm, overall 6.86 mm). Body 6.10 × 2.69 mm. Courtyard 7.4 × 3.7 mm (Littelfuse 451 ``Recommended pad layout'') \\
SMF24A (unidirectional, SOD-123FL) & Stock \texttt{Device:\allowbreak{}D\_\allowbreak{}Zener} (pin 1 K, 2 A), Value = \texttt{SMF24A}. Don't use \texttt{Diode:\allowbreak{}SMF24A}; its A1/A2 pin names are ambiguous & Stock \texttt{Diode\_\allowbreak{}SMD:\allowbreak{}D\_\allowbreak{}SOD-123F} (pad 1 = cathode band) \\
\end{longtable}
}

Everything else uses stock symbols/footprints, named per row below.

\paragraph{0.1 Custom symbol pin tables (electrical types for ERC)}

\texttt{power:\allowbreak{}TPS61178}: 1 \texttt{FREQ/\allowbreak{}SYNC} input · 2 \texttt{AGND} power\_\allowbreak{}in · 3 \texttt{ILIMIT} passive · 4 \texttt{COMP} passive · 5 \texttt{FB} input · 6 \texttt{PGND} power\_\allowbreak{}in · 7 \texttt{SW} power\_\allowbreak{}out · 8 \texttt{VOUT} power\_\allowbreak{}out · 9 \texttt{DISDRV} output · 10 \texttt{VCC} power\_\allowbreak{}out · 11 \texttt{EN} input · 12 \texttt{VIN} power\_\allowbreak{}in · 13 \texttt{BST} passive.

\texttt{power:\allowbreak{}CH224Q}: 1 \texttt{VHV} power\_\allowbreak{}in · 2 \texttt{CFG2/\allowbreak{}SCL} bidirectional · 3 \texttt{CFG3/\allowbreak{}SDA} bidirectional · 4 \texttt{DP} bidirectional · 5 \texttt{DM} bidirectional · 6 \texttt{CC2} bidirectional · 7 \texttt{CC1} bidirectional · 8 \texttt{VBUS} input · 9 \texttt{CFG1} input · 10 \texttt{PG} open\_\allowbreak{}collector · 11 \texttt{GND} power\_\allowbreak{}in (exposed pad, ``pin 0'' in the WCH table).

\texttt{power:\allowbreak{}R5460x2}: 1 \texttt{DOUT} output · 2 \texttt{COUT} output · 3 \texttt{V-} input · 4 \texttt{VC} input · 5 \texttt{VDD} power\_\allowbreak{}in · 6 \texttt{VSS} power\_\allowbreak{}in.

\texttt{power:\allowbreak{}HY2213}: 1 \texttt{NC} no\_\allowbreak{}connect · 2 \texttt{VDD} power\_\allowbreak{}in · 3 \texttt{VSS} power\_\allowbreak{}in · 4 \texttt{NC} no\_\allowbreak{}connect · 5 \texttt{NC} no\_\allowbreak{}connect · 6 \texttt{OUT} output.

\begin{center}\rule{0.5\linewidth}{0.5pt}\end{center}

\subsubsection{Net list summary}

{\def\LTcaptype{none} % do not increment counter
\begin{longtable}[]{@{}
  >{\raggedright\arraybackslash}p{(\linewidth - 4\tabcolsep) * \real{0.3333}}
  >{\raggedright\arraybackslash}p{(\linewidth - 4\tabcolsep) * \real{0.3333}}
  >{\raggedright\arraybackslash}p{(\linewidth - 4\tabcolsep) * \real{0.3333}}@{}}
\toprule\noalign{}
\begin{minipage}[b]{\linewidth}\raggedright
Net
\end{minipage} & \begin{minipage}[b]{\linewidth}\raggedright
Scope
\end{minipage} & \begin{minipage}[b]{\linewidth}\raggedright
Description / producer → consumers
\end{minipage} \\
\midrule\noalign{}
\endhead
\bottomrule\noalign{}
\endlastfoot
\texttt{VBUS} & \textbf{global} & USB-C VBUS (5 V at attach, 15 V after PD). J601 → U603, U701 VBUS/VAC1/VAC2, D601 \\
\texttt{VSYS} & \textbf{global} & BQ25792 SYS output, 6.0--8.6 V (≥ 7.0 V with VBUS present). U701 → U901 boost, U951 buck \\
\texttt{VBAT\_\allowbreak{}PACK} & \textbf{global} & Pack + after F801 fuse. Battery sheet → Q701 drain, U701 BATP (via R706) \\
\texttt{PVDD} & \textbf{global} & Amp rail, 18.0 V = TPS61178 VOUT (no disconnect FET) → amp sheets. ≈ VSYS − 0.5 V when BOOST\_\allowbreak{}EN is low \\
\texttt{+3V3} & \textbf{global} & Buck output → everything logic \\
\texttt{GND} & \textbf{global} & System ground = protected pack negative (P−) \\
\texttt{I2C\_\allowbreak{}SDA}, \texttt{I2C\_\allowbreak{}SCL} & \textbf{global} & Shared bus. Pull-ups 2.2 kΩ on mcu sheet (R104/R105) \\
\texttt{USB\_\allowbreak{}DP}, \texttt{USB\_\allowbreak{}DN} & \textbf{global} & USB-C D+/D− → CH340C on \texttt{usb\_\allowbreak{}uart} \\
\texttt{CHG\_\allowbreak{}INT} & \textbf{global} & BQ25792 /INT (open drain, 256 µs low pulse). 10 k pull-up on mcu sheet (R106) \\
\texttt{BOOST\_\allowbreak{}EN} & \textbf{global} & ESP32 GPIO2 (strap) → TPS61178 EN. 100 k pull-down on mcu sheet (R107). No pull-up anywhere; the TPS61178 EN has only an internal 800 kΩ pull-down \\
\texttt{PVDD\_\allowbreak{}TRIM} & \textbf{global} & ESP32 GPIO25 (DAC1 or LEDC PWM) → trim RC on rails sheet \\
\texttt{BAT\_\allowbreak{}TS} & \textbf{global (new)} & BQ25792 TS node: RT1/RT2 on charger sheet, NTC select on battery sheet \\
\texttt{CHG\_\allowbreak{}QON} & \textbf{global} & BQ25792 /QON (internal pull-up). SW701 on the main board \textbf{and} the UI-board power button through FFC J501: one former GND pin, 19 or 20, as assigned in spec\_\allowbreak{}signal §4.2 (ECO-002 \#2). Button to GND only, ESD on the UI board, no pull-up and no driven signal \\
\texttt{CHG\_\allowbreak{}CE} & not used & /CE is tied low through R707 (charging allowed by hardware; firmware uses EN\_\allowbreak{}CHG). No GPIO needed \\
\end{longtable}
}

Local nets are defined per sheet.

\begin{center}\rule{0.5\linewidth}{0.5pt}\end{center}

\subsubsection{\texorpdfstring{1. Sheet \texttt{usb\_\allowbreak{}power}}{1. Sheet usb\_\allowbreak{}power}}

\paragraph{1.1 Component table}

{\def\LTcaptype{none} % do not increment counter
\begin{longtable}[]{@{}
  >{\raggedright\arraybackslash}p{(\linewidth - 14\tabcolsep) * \real{0.1250}}
  >{\raggedright\arraybackslash}p{(\linewidth - 14\tabcolsep) * \real{0.1250}}
  >{\raggedright\arraybackslash}p{(\linewidth - 14\tabcolsep) * \real{0.1250}}
  >{\raggedright\arraybackslash}p{(\linewidth - 14\tabcolsep) * \real{0.1250}}
  >{\raggedright\arraybackslash}p{(\linewidth - 14\tabcolsep) * \real{0.1250}}
  >{\raggedright\arraybackslash}p{(\linewidth - 14\tabcolsep) * \real{0.1250}}
  >{\raggedright\arraybackslash}p{(\linewidth - 14\tabcolsep) * \real{0.1250}}
  >{\raggedright\arraybackslash}p{(\linewidth - 14\tabcolsep) * \real{0.1250}}@{}}
\toprule\noalign{}
\begin{minipage}[b]{\linewidth}\raggedright
Ref
\end{minipage} & \begin{minipage}[b]{\linewidth}\raggedright
Value
\end{minipage} & \begin{minipage}[b]{\linewidth}\raggedright
Footprint
\end{minipage} & \begin{minipage}[b]{\linewidth}\raggedright
Rating
\end{minipage} & \begin{minipage}[b]{\linewidth}\raggedright
LCSC
\end{minipage} & \begin{minipage}[b]{\linewidth}\raggedright
B/E
\end{minipage} & \begin{minipage}[b]{\linewidth}\raggedright
Symbol
\end{minipage} & \begin{minipage}[b]{\linewidth}\raggedright
Notes
\end{minipage} \\
\midrule\noalign{}
\endhead
\bottomrule\noalign{}
\endlastfoot
J601 & USB-C 16P TYPE-C-31-M-12 & \texttt{Connector\_\allowbreak{}USB:\allowbreak{}USB\_\allowbreak{}C\_\allowbreak{}Receptacle\_\allowbreak{}HRO\_\allowbreak{}TYPE-C-31-M-12} & 20 V 5 A & C165948 & E & \texttt{Connector:\allowbreak{}USB\_\allowbreak{}C\_\allowbreak{}Receptacle\_\allowbreak{}USB2.0\_\allowbreak{}16P} & stock sym+fp \\
D601 & SMF24A & \texttt{Diode\_\allowbreak{}SMD:\allowbreak{}D\_\allowbreak{}SOD-123F} & VRWM 24 V, VBR 26.7--29.5 V, Vc 38.9 V @ 5.1 A, 200 W & C123810 & E & \texttt{Device:\allowbreak{}D\_\allowbreak{}Zener} & VBUS TVS at connector \\
C601 & 100 nF & C\_\allowbreak{}0603 & 50 V X7R & C14663 & B & \texttt{Device:\allowbreak{}C} & VBUS HF at connector \\
R601 & 2.0 Ω & R\_\allowbreak{}0805 & 0.125 W & C2930265 & E & \texttt{Device:\allowbreak{}R} & \textbf{DNP}. Hot-plug snubber (BQ25798 Fig 8-1 note 1) \\
C602 & 2.2 µF & C\_\allowbreak{}0805 & 50 V X5R & C377773 & B & \texttt{Device:\allowbreak{}C} & \textbf{DNP}, in series with R601 \\
U601 & USBLC6-2SC6 & \texttt{Package\_\allowbreak{}TO\_\allowbreak{}SOT\_\allowbreak{}SMD:\allowbreak{}SOT-23-6} & 5 V, 6 A 8/20 µs & C2687116 & E & \texttt{Power\_\allowbreak{}Protection:\allowbreak{}USBLC6-2SC6} & D+/D− ESD \\
C603 & 100 nF & C\_\allowbreak{}0603 & 50 V & C14663 & B & \texttt{Device:\allowbreak{}C} & U601 pin 5 to GND \\
U602 & USBLC6-2SC6 & SOT-23-6 & & C2687116 & E & \texttt{Power\_\allowbreak{}Protection:\allowbreak{}USBLC6-2SC6} & CC1/CC2 ESD (pin 5 not connected) \\
U603 & CH224Q & \texttt{Package\_\allowbreak{}DFN\_\allowbreak{}QFN:\allowbreak{}Texas\_\allowbreak{}DSQ0010A\_\allowbreak{}WSON-10-1EP\_\allowbreak{}2x2mm\_\allowbreak{}P0.4mm\_\allowbreak{}EP0.9x1.5mm} & VHV ≤ 32 V abs & C46061833 & E & \textbf{custom} \texttt{power:\allowbreak{}CH224Q} & PD sink \\
C604 & 1 µF & C\_\allowbreak{}0603 & \textbf{50 V} X5R & C15849 & B & \texttt{Device:\allowbreak{}C} & VHV decoupling (``注意耐压'': mind the voltage rating) \\
R602 & 56 kΩ 1 \% & R\_\allowbreak{}0603 & & C23206 & B & \texttt{Device:\allowbreak{}R} & CFG1 → GND: single-resistor mode, \textbf{15 V default}, enables I2C (CH224 Table 5-1, §5.2.3) \\
Q601, Q602 & AO3400A & \texttt{Package\_\allowbreak{}TO\_\allowbreak{}SOT\_\allowbreak{}SMD:\allowbreak{}SOT-23} & 30 V, Vgs ±12 V, Vth 0.65--1.45 V & C20917 & B & \texttt{Transistor\_\allowbreak{}FET:\allowbreak{}AO3400A} (1 G, 2 S, 3 D) & I2C isolators for CH224Q \\
R603 & 100 kΩ & R\_\allowbreak{}0603 & & C25803 & B & \texttt{Device:\allowbreak{}R} & VBUS → \texttt{PD\_\allowbreak{}ISO\_\allowbreak{}G} \\
R604 & 100 kΩ & R\_\allowbreak{}0603 & & C25803 & B & \texttt{Device:\allowbreak{}R} & \texttt{PD\_\allowbreak{}ISO\_\allowbreak{}G} → GND \\
TP601 & VBUS & \texttt{TestPoint:\allowbreak{}TestPoint\_\allowbreak{}Pad\_\allowbreak{}D1.5mm} & & -- & -- & \texttt{Connector:\allowbreak{}TestPoint} & \\
TP602 & PD\_\allowbreak{}PG & \texttt{TestPoint:\allowbreak{}TestPoint\_\allowbreak{}Pad\_\allowbreak{}D1.5mm} & & -- & -- & \texttt{Connector:\allowbreak{}TestPoint} & optional \\
\end{longtable}
}

\paragraph{1.2 Pin-by-pin connections}

\textbf{J601 USB-C} (\texttt{USB\_\allowbreak{}C\_\allowbreak{}Receptacle\_\allowbreak{}USB2.0\_\allowbreak{}16P}):

{\def\LTcaptype{none} % do not increment counter
\begin{longtable}[]{@{}lll@{}}
\toprule\noalign{}
Pin & Name & Net \\
\midrule\noalign{}
\endhead
\bottomrule\noalign{}
\endlastfoot
A1, A12, B1, B12 & GND & \texttt{GND} \\
A4, A9, B4, B9 & VBUS & \texttt{VBUS} \\
A5 & CC1 & \texttt{CC1} \\
B5 & CC2 & \texttt{CC2} \\
A6, B6 & D+ & \texttt{USB\_\allowbreak{}DP} (global) \\
A7, B7 & D− & \texttt{USB\_\allowbreak{}DN} (global) \\
A8 & SBU1 & no connect \\
B8 & SBU2 & no connect \\
SH & SHIELD & \texttt{GND} \\
\end{longtable}
}

\textbf{D601 SMF24A}: pin 1 (K) \texttt{VBUS}, pin 2 (A) \texttt{GND}. Place within 3 mm of J601's VBUS pins.

\textbf{R601/C602 snubber (DNP)}: R601 pin 1 \texttt{VBUS}, pin 2 \texttt{VBUS\_\allowbreak{}SNUB}; C602 pin 1 \texttt{VBUS\_\allowbreak{}SNUB}, pin 2 \texttt{GND}.

\textbf{U601 USBLC6-2SC6 (D+/D−)}:

{\def\LTcaptype{none} % do not increment counter
\begin{longtable}[]{@{}
  >{\raggedright\arraybackslash}p{(\linewidth - 4\tabcolsep) * \real{0.3333}}
  >{\raggedright\arraybackslash}p{(\linewidth - 4\tabcolsep) * \real{0.3333}}
  >{\raggedright\arraybackslash}p{(\linewidth - 4\tabcolsep) * \real{0.3333}}@{}}
\toprule\noalign{}
\begin{minipage}[b]{\linewidth}\raggedright
Pin
\end{minipage} & \begin{minipage}[b]{\linewidth}\raggedright
Name
\end{minipage} & \begin{minipage}[b]{\linewidth}\raggedright
Net
\end{minipage} \\
\midrule\noalign{}
\endhead
\bottomrule\noalign{}
\endlastfoot
1 & I/O1 & \texttt{USB\_\allowbreak{}DP} \\
2 & GND & \texttt{GND} \\
3 & I/O2 & \texttt{USB\_\allowbreak{}DN} \\
4 & I/O2 & \texttt{USB\_\allowbreak{}DN} \\
5 & VBUS & \texttt{+3V3} (\textbf{not} VBUS: VBUS reaches 15--20 V and the USBLC6 rail clamps at 6--9 V). C603 to GND \\
6 & I/O1 & \texttt{USB\_\allowbreak{}DP} \\
\end{longtable}
}

\textbf{U602 USBLC6-2SC6 (CC)}: 1 \texttt{CC1}, 2 \texttt{GND}, 3 \texttt{CC2}, 4 \texttt{CC2}, 5 \textbf{no connect} (the CC pins can see Vconn up to 5.5 V, so they must not be steered into +3V3; VBR(min) 6 V \textgreater{} 5.5 V), 6 \texttt{CC1}.

\textbf{U603 CH224Q}:

{\def\LTcaptype{none} % do not increment counter
\begin{longtable}[]{@{}
  >{\raggedright\arraybackslash}p{(\linewidth - 6\tabcolsep) * \real{0.2500}}
  >{\raggedright\arraybackslash}p{(\linewidth - 6\tabcolsep) * \real{0.2500}}
  >{\raggedright\arraybackslash}p{(\linewidth - 6\tabcolsep) * \real{0.2500}}
  >{\raggedright\arraybackslash}p{(\linewidth - 6\tabcolsep) * \real{0.2500}}@{}}
\toprule\noalign{}
\begin{minipage}[b]{\linewidth}\raggedright
Pin
\end{minipage} & \begin{minipage}[b]{\linewidth}\raggedright
Name
\end{minipage} & \begin{minipage}[b]{\linewidth}\raggedright
Net
\end{minipage} & \begin{minipage}[b]{\linewidth}\raggedright
Notes
\end{minipage} \\
\midrule\noalign{}
\endhead
\bottomrule\noalign{}
\endlastfoot
1 & VHV & \texttt{VBUS} & C604 1 µF/50 V to GND at the pin \\
2 & CFG2/SCL & \texttt{PD\_\allowbreak{}SCL} & internal 10 kΩ pull-up (7--15 k) \\
3 & CFG3/SDA & \texttt{PD\_\allowbreak{}SDA} & internal 10 kΩ pull-up \\
4 & DP & no connect & USB data goes only to the CH340C. BC1.2/QC on the CH224Q is unused \\
5 & DM & no connect & \\
6 & CC2 & \texttt{CC2} & Rd is internal (no external 5.1 k, per CH224 §6.1.1) \\
7 & CC1 & \texttt{CC1} & \\
8 & VBUS & \texttt{VBUS} & ``需要短接至 VHV'': must be shorted to VHV \\
9 & CFG1 & \texttt{PD\_\allowbreak{}CFG1} & R602 56 k to GND \\
10 & PG & \texttt{PD\_\allowbreak{}PG} & open drain, low = PD power good. TP602 only \\
11 (EP) & GND & \texttt{GND} & EP to GND with ≥ 2 vias \\
\end{longtable}
}

\textbf{Q601 / Q602 I2C isolator} (source on the CH224Q side, drain on the bus side):

{\def\LTcaptype{none} % do not increment counter
\begin{longtable}[]{@{}llll@{}}
\toprule\noalign{}
Ref & G (1) & S (2) & D (3) \\
\midrule\noalign{}
\endhead
\bottomrule\noalign{}
\endlastfoot
Q601 & \texttt{PD\_\allowbreak{}ISO\_\allowbreak{}G} & \texttt{PD\_\allowbreak{}SCL} & \texttt{I2C\_\allowbreak{}SCL} \\
Q602 & \texttt{PD\_\allowbreak{}ISO\_\allowbreak{}G} & \texttt{PD\_\allowbreak{}SDA} & \texttt{I2C\_\allowbreak{}SDA} \\
\end{longtable}
}

R603: \texttt{VBUS} ↔ \texttt{PD\_\allowbreak{}ISO\_\allowbreak{}G}. R604: \texttt{PD\_\allowbreak{}ISO\_\allowbreak{}G} ↔ \texttt{GND}.

\paragraph{1.3 Design notes / calculations}

\begin{itemize}
\tightlist
\item
  \textbf{PD request.} CFG1 = 56 kΩ → 15 V at attach, and I2C is enabled automatically in single-resistor mode (CH224 §5.2.1, §5.2.3). Firmware can rewrite reg 0x0A (0 = 5 V, 1 = 9, 2 = 12, 3 = 15, 4 = 20 V) and read reg 0x50 (max current in 50 mA units) to program the BQ IINDPM. \textbf{Never request 28 V/EPR}: the BQ25792 VBUS operating max is 24 V, and the isolator gate would see 14 V (\textgreater{} AO3400A's ±12 V Vgs).
\item
  \textbf{CH224Q I2C address} is ``0x22 or 0x23'' (CH224 §5.2.3; the manual doesn't say which). Firmware should probe both. Bus map: CH224Q 0x22/0x23, TAS5805M 0x2C (4.7 kΩ ADR) and 0x2D (15 kΩ ADR) (TAS5805M Table 7-5), BQ25792 0x6B (BQ25798 §7.3.14.5). \textbf{No collisions.}
\item
  \textbf{Isolator.} Gate = VBUS/2: 2.5 V at 5 V and 7.5 V at 15 V; 10 V at 20 V (within the ±12 V Vgs rating).

  \begin{itemize}
  \tightlist
  \item
    VBUS absent: gate 0 V, so Vgs = 0 and the FET is off. The body diode (anode = source = CH224Q side ≈ 0 V, cathode = drain = bus at 3.3 V) is reverse-biased, so the bus is isolated.
  \item
    VBUS present, CH224Q pulls low: Vgs ≥ 2.5 V \textgreater{} Vth(max) 1.45 V, so the FET is on.
  \item
    VBUS present, ESP32 pulls low: the body diode drags the source to ≈0.6 V, so Vgs ≥ 1.9 V and the FET turns on. This is the standard bidirectional NMOS level shifter.
  \item
    Added bus capacitance is ≈ 2 × 75 pF Coss. That fits within the 400 pF budget at 2.2 kΩ pull-ups (spec\_\allowbreak{}signal §I2C).
  \end{itemize}
\item
  \textbf{Back-powering summary (brief question).}

  \begin{itemize}
  \tightlist
  \item
    BQ25792: powered from BAT (or VBUS) whenever +3V3 exists, since +3V3 derives from VSYS. Its I2C pins are open-drain and rated 6 V abs max. No risk.
  \item
    TAS5805M: DVDD = +3V3, the same rail as the pull-ups. No risk.
  \item
    CH224Q: powered only from VBUS. \textbf{Risk, handled by Q601/Q602.} Its CFG2/CFG3 pins tolerate 6.5 V (CH224 §8.1.1), and V\_\allowbreak{}IH(I2C) min is 1.5 V, so 3.3 V logic is fine when it is powered.
  \end{itemize}
\item
  \textbf{ESD.} VBUS gets D601 (SMF24A). CC and D± get USBLC6. USB-C sources always attach at 5 V and only step to 15 V after the PD contract, so the BQ ``\textgreater15 V hot-plug'' snubber (R601/C602) stays DNP.
\item
  \textbf{VBUS capacitance.} Total sink capacitance (VBUS + PMID, which are DC-connected inside the BQ through its input current sense) is about 50 µF nominal, ≈ 20 µF effective at 15 V. That exceeds the USB Type-C cSnkBulk guideline of 10 µF. TI's recommended values and EVM do the same. Accepted: the source limits inrush at 5 V attach.
\end{itemize}

\begin{center}\rule{0.5\linewidth}{0.5pt}\end{center}

\subsubsection{\texorpdfstring{2. Sheet \texttt{charger} (BQ25792 NVDC buck-boost charger)}{2. Sheet charger (BQ25792 NVDC buck-boost charger)}}

\paragraph{2.1 Component table}

{\def\LTcaptype{none} % do not increment counter
\begin{longtable}[]{@{}
  >{\raggedright\arraybackslash}p{(\linewidth - 14\tabcolsep) * \real{0.1250}}
  >{\raggedright\arraybackslash}p{(\linewidth - 14\tabcolsep) * \real{0.1250}}
  >{\raggedright\arraybackslash}p{(\linewidth - 14\tabcolsep) * \real{0.1250}}
  >{\raggedright\arraybackslash}p{(\linewidth - 14\tabcolsep) * \real{0.1250}}
  >{\raggedright\arraybackslash}p{(\linewidth - 14\tabcolsep) * \real{0.1250}}
  >{\raggedright\arraybackslash}p{(\linewidth - 14\tabcolsep) * \real{0.1250}}
  >{\raggedright\arraybackslash}p{(\linewidth - 14\tabcolsep) * \real{0.1250}}
  >{\raggedright\arraybackslash}p{(\linewidth - 14\tabcolsep) * \real{0.1250}}@{}}
\toprule\noalign{}
\begin{minipage}[b]{\linewidth}\raggedright
Ref
\end{minipage} & \begin{minipage}[b]{\linewidth}\raggedright
Value
\end{minipage} & \begin{minipage}[b]{\linewidth}\raggedright
Footprint
\end{minipage} & \begin{minipage}[b]{\linewidth}\raggedright
Rating
\end{minipage} & \begin{minipage}[b]{\linewidth}\raggedright
LCSC
\end{minipage} & \begin{minipage}[b]{\linewidth}\raggedright
B/E
\end{minipage} & \begin{minipage}[b]{\linewidth}\raggedright
Symbol
\end{minipage} & \begin{minipage}[b]{\linewidth}\raggedright
Notes / citation
\end{minipage} \\
\midrule\noalign{}
\endhead
\bottomrule\noalign{}
\endlastfoot
U701 & BQ25792RQMR & \texttt{Package\_\allowbreak{}DFN\_\allowbreak{}QFN:\allowbreak{}Texas\_\allowbreak{}RQM0029A\_\allowbreak{}VQFN-29\_\allowbreak{}4x4mm\_\allowbreak{}P0.4mm} & & C2862876 & E & \texttt{Battery\_\allowbreak{}Management:\allowbreak{}BQ25798} (Value BQ25792RQMR) & pinout identical (§0) \\
L701 & 2.2 µH FXL0630-2R2-M & \texttt{Inductor\_\allowbreak{}SMD:\allowbreak{}L\_\allowbreak{}Changjiang\_\allowbreak{}FXL0630} & Isat 10 A / Irms 9.5 A, 15 mΩ, 7×6.6 mm & C167218 & E & \texttt{Device:\allowbreak{}L} & SW1--SW2. Alt: FXL1040-2R2-M C167228 (same as boost, saves a line item, 11.5×10 mm) \\
C701, C702 & 10 µF & C\_\allowbreak{}1206 & \textbf{50 V} X5R & C13585 & B & \texttt{Device:\allowbreak{}C} & VBUS (2×10 µF + 0.1 µF, Table 5-1) \\
C703 & 100 nF & C\_\allowbreak{}0603 & 50 V X7R & C14663 & B & \texttt{Device:\allowbreak{}C} & VBUS, closest to pins 2/3 \\
C704, C705, C706 & 10 µF & C\_\allowbreak{}1206 & \textbf{50 V} X5R & C13585 & B & \texttt{Device:\allowbreak{}C} & PMID (3×10 µF + 0.1 µF, §8.2.2.3) \\
C707 & 100 nF & C\_\allowbreak{}0603 & 50 V & C14663 & B & \texttt{Device:\allowbreak{}C} & PMID, closest to pin 29 \\
C708 & 47 nF & C\_\allowbreak{}0603 & 50 V X7R & C1622 & B & \texttt{Device:\allowbreak{}C} & BTST1--SW1 (≥10 V, 47 nF, Table 5-1) \\
C709 & 47 nF & C\_\allowbreak{}0603 & 50 V X7R & C1622 & B & \texttt{Device:\allowbreak{}C} & BTST2--SW2 \\
C710 & 4.7 µF & C\_\allowbreak{}0603 & 16 V X5R & C19666 & B & \texttt{Device:\allowbreak{}C} & REGN (10 V+ rated, 4.7 µF) \\
C711--C715 & 10 µF & C\_\allowbreak{}0805 & 25 V X5R & C15850 & B & \texttt{Device:\allowbreak{}C} & SYS (5×10 µF + 0.1 µF, §8.2.2.4; 16 V+ for 2S) \\
C716 & 100 nF & C\_\allowbreak{}0603 & 50 V & C14663 & B & \texttt{Device:\allowbreak{}C} & SYS, closest to pin 25 \\
C717, C718 & 10 µF & C\_\allowbreak{}0805 & 25 V X5R & C15850 & B & \texttt{Device:\allowbreak{}C} & BAT (2×10 µF, Table 5-1) \\
C719 & 47 nF & C\_\allowbreak{}0603 & 50 V & C1622 & B & \texttt{Device:\allowbreak{}C} & \textbf{DNP}. BATP filter (Fig 8-2 note 2, 4S only) \\
C720 & 10 nF & C\_\allowbreak{}0603 & 50 V X7R & C57112 & B & \texttt{Device:\allowbreak{}C} & \textbf{DNP}. Ship-FET gate--\textbf{drain} (\texttt{SFET\_\allowbreak{}G}--\texttt{VBAT\_\allowbreak{}PACK}). Slows the source-follower ramp at battery connect. \textbf{Fitting rule: fit C720 only together with C721 ≥ 5 × C720, or with R710} \\
C721 & 100 nF & C\_\allowbreak{}0603 & 50 V X7R & C14663 & B & \texttt{Device:\allowbreak{}C} & \textbf{DNP} (ECO-002 \#1, S2). Ship-FET gate--\textbf{source} (\texttt{SFET\_\allowbreak{}G}--\texttt{CHG\_\allowbreak{}BAT}), placed next to C720. Holds the gate down against C720's coupling at pack hot-plug \\
R701 & 8.2 kΩ 1 \% & R\_\allowbreak{}0603 & & C25981 & B & \texttt{Device:\allowbreak{}R} & PROG → GND: \textbf{2S, 750 kHz} (Table 7-1; EVM R30 8.20 k) \\
R702 & 51 kΩ 1 \% & R\_\allowbreak{}0603 & & C23196 & B & \texttt{Device:\allowbreak{}R} & REGN → ILIM\_\allowbreak{}HIZ \\
R703 & 100 kΩ 1 \% & R\_\allowbreak{}0603 & & C25803 & B & \texttt{Device:\allowbreak{}R} & ILIM\_\allowbreak{}HIZ → GND \\
R704 & 5.1 kΩ 1 \% & R\_\allowbreak{}0603 & & C23186 & B & \texttt{Device:\allowbreak{}R} & RT1: REGN → BAT\_\allowbreak{}TS \\
R705 & 30 kΩ 1 \% & R\_\allowbreak{}0603 & & C22984 & B & \texttt{Device:\allowbreak{}R} & RT2: BAT\_\allowbreak{}TS → GND \\
R706 & 100 Ω & R\_\allowbreak{}0603 & & C22775 & B & \texttt{Device:\allowbreak{}R} & VBAT\_\allowbreak{}PACK → BATP (Table 5-1: ``100 Ω series'') \\
R707 & 10 kΩ & R\_\allowbreak{}0603 & & C25804 & B & \texttt{Device:\allowbreak{}R} & /CE → GND (charge enabled in hardware) \\
R708 & 1.0 kΩ & R\_\allowbreak{}0603 & & C21190 & B & \texttt{Device:\allowbreak{}R} & STAT LED series \\
D701 & Green LED KT-0805G & \texttt{LED\_\allowbreak{}SMD:\allowbreak{}LED\_\allowbreak{}0805\_\allowbreak{}2012Metric} & Vf 2.6--3.1 V & C2297 & B & \texttt{Device:\allowbreak{}LED} & Charge status \\
SW701 & TS-1187A-B-A-B & \texttt{Button\_\allowbreak{}Switch\_\allowbreak{}SMD:\allowbreak{}SW\_\allowbreak{}Push\_\allowbreak{}1P1T\_\allowbreak{}XKB\_\allowbreak{}TS-1187A} & 12 V 50 mA & C318884 & B & \texttt{Switch:\allowbreak{}SW\_\allowbreak{}Push} & QON: wake from ship mode / 10 s = system reset. Paralleled by the UI-board button through the FFC \\
Q701 & NCEP3065QU & \texttt{Package\_\allowbreak{}SON:\allowbreak{}VSON-8\_\allowbreak{}3.3x3.3mm\_\allowbreak{}P0.65mm\_\allowbreak{}NexFET} & 30 V, 1.9 mΩ@10 V, 3.0 mΩ@4.5 V, Vth 1.0--2.0 V & C502964 & E (same part as Q801/Q802) & \texttt{Transistor\_\allowbreak{}FET:\allowbreak{}CSD17581Q3A} (Value NCEP3065QU) & \textbf{Ship FET (SFET)} \\
R709 & 0 Ω & R\_\allowbreak{}0603 & & C21189 & B & \texttt{Device:\allowbreak{}R} & SDRV → SFET gate (EVM R17) \\
R710 & 100 kΩ & R\_\allowbreak{}0603 & & C25803 & B & \texttt{Device:\allowbreak{}R} & \textbf{DNP}. SFET gate--source (EVM R16) \\
TP701 & VSYS & \texttt{TestPoint:\allowbreak{}TestPoint\_\allowbreak{}Pad\_\allowbreak{}D1.5mm} & & -- & -- & \texttt{Connector:\allowbreak{}TestPoint} & \\
TP702 & CHG\_\allowbreak{}BAT & \texttt{TestPoint:\allowbreak{}TestPoint\_\allowbreak{}Pad\_\allowbreak{}D1.5mm} & & -- & -- & \texttt{Connector:\allowbreak{}TestPoint} & optional \\
\end{longtable}
}

\paragraph{2.2 U701 pin-by-pin}

{\def\LTcaptype{none} % do not increment counter
\begin{longtable}[]{@{}
  >{\raggedright\arraybackslash}p{(\linewidth - 6\tabcolsep) * \real{0.2500}}
  >{\raggedright\arraybackslash}p{(\linewidth - 6\tabcolsep) * \real{0.2500}}
  >{\raggedright\arraybackslash}p{(\linewidth - 6\tabcolsep) * \real{0.2500}}
  >{\raggedright\arraybackslash}p{(\linewidth - 6\tabcolsep) * \real{0.2500}}@{}}
\toprule\noalign{}
\begin{minipage}[b]{\linewidth}\raggedright
Pin
\end{minipage} & \begin{minipage}[b]{\linewidth}\raggedright
Name
\end{minipage} & \begin{minipage}[b]{\linewidth}\raggedright
Net
\end{minipage} & \begin{minipage}[b]{\linewidth}\raggedright
Connection
\end{minipage} \\
\midrule\noalign{}
\endhead
\bottomrule\noalign{}
\endlastfoot
1 & STAT & \texttt{CHG\_\allowbreak{}STAT} & D701 cathode. D701 anode → R708 → \texttt{CHG\_\allowbreak{}REGN} (Fig 8-1: REGN--2.2 k--LED--STAT) \\
2, 3 & VBUS & \texttt{VBUS} & C701, C702, C703 \\
4 & BTST1 & \texttt{CHG\_\allowbreak{}BTST1} & C708 to \texttt{CHG\_\allowbreak{}SW1} \\
5 & REGN & \texttt{CHG\_\allowbreak{}REGN} & C710 to GND; feeds R702, R704, R708 only \\
6 & D+ & no connect & (see §2.3, input current) \\
7 & D− & no connect & \\
8 & VAC2 & \texttt{VBUS} & ``Connect to VBUS if ACFET2/RBFET2 not installed'' \\
9 & VAC1 & \texttt{VBUS} & same, for port 1 \\
10 & ACDRV2 & \texttt{GND} & ``Tie to GND if no ACFET2/RBFET2'' \\
11 & ACDRV1 & \texttt{GND} & ``Tie to GND if no ACFET1/RBFET1'' \\
12 & /QON & \texttt{CHG\_\allowbreak{}QON} (global) & SW701 to GND, plus the UI button through J501. Internal pull-up (≈3.6--3.8 V), no external resistor \\
13 & /CE & \texttt{CHG\_\allowbreak{}CE\_\allowbreak{}N} & R707 10 k to GND (must not float) \\
14 & SCL & \texttt{I2C\_\allowbreak{}SCL} & pull-up on mcu sheet \\
15 & SDA & \texttt{I2C\_\allowbreak{}SDA} & pull-up on mcu sheet \\
16 & TS & \texttt{BAT\_\allowbreak{}TS} (global) & R704 from REGN, R705 to GND, NTC on battery sheet \\
17 & ILIM\_\allowbreak{}HIZ & \texttt{CHG\_\allowbreak{}ILIM} & R702 from REGN, R703 to GND \\
18 & BATP & \texttt{CHG\_\allowbreak{}BATP} & R706 from \texttt{VBAT\_\allowbreak{}PACK} (battery side of the ship FET, EVM JP9 ``BATTERY''); C719 DNP to GND \\
19 & BTST2 & \texttt{CHG\_\allowbreak{}BTST2} & C709 to \texttt{CHG\_\allowbreak{}SW2} \\
20 & PROG & \texttt{CHG\_\allowbreak{}PROG} & R701 8.2 k to GND \\
21 & /INT & \texttt{CHG\_\allowbreak{}INT} (global) & pull-up on mcu sheet (to +3V3, as §7.3.4.1 recommends, not to REGN) \\
22, 23 & BAT & \texttt{CHG\_\allowbreak{}BAT} & C717, C718; Q701 source \\
24 & SDRV & \texttt{CHG\_\allowbreak{}SDRV} & R709 0 Ω → \texttt{SFET\_\allowbreak{}G} \\
25 & SYS & \texttt{VSYS} & C711--C716 \\
26 & SW2 & \texttt{CHG\_\allowbreak{}SW2} & L701 pin 2 \\
27 & GND & \texttt{GND} & power ground, via array \\
28 & SW1 & \texttt{CHG\_\allowbreak{}SW1} & L701 pin 1 \\
29 & PMID & \texttt{CHG\_\allowbreak{}PMID} & C704--C707 only (no load) \\
\end{longtable}
}

\textbf{L701}: pin 1 \texttt{CHG\_\allowbreak{}SW1}, pin 2 \texttt{CHG\_\allowbreak{}SW2}.

\textbf{Q701 ship FET} (orientation as TI EVM Q5: source on the charger BAT side, drain on the battery side, so the body diode blocks battery → system when off):

{\def\LTcaptype{none} % do not increment counter
\begin{longtable}[]{@{}ll@{}}
\toprule\noalign{}
Pin & Net \\
\midrule\noalign{}
\endhead
\bottomrule\noalign{}
\endlastfoot
1, 2, 3 (S) & \texttt{CHG\_\allowbreak{}BAT} \\
4 (G) & \texttt{SFET\_\allowbreak{}G} \\
5 (D, tab) & \texttt{VBAT\_\allowbreak{}PACK} \\
\end{longtable}
}

R710 (DNP): \texttt{SFET\_\allowbreak{}G} ↔ \texttt{CHG\_\allowbreak{}BAT}. C720 (DNP): \texttt{SFET\_\allowbreak{}G} ↔ \texttt{VBAT\_\allowbreak{}PACK}. C721 (DNP): \texttt{SFET\_\allowbreak{}G} ↔ \texttt{CHG\_\allowbreak{}BAT}.

\textbf{Why the C720 rule.} C720 alone is unsafe. If SDRV does not actively hold the gate low while U701 is still unpowered (the datasheet doesn't say), a pack hot-plug couples the gate to 8.4 V × 10 nF/(10 nF + 2.1 nF Ciss) ≈ 6.9 V. Q701 would then turn fully on into ≈1.5 mF and could trip the R5460.

\begin{itemize}
\tightlist
\item
  With C721 = 100 nF: 8.4 × 10/(10 + 100 + 2.1) ≈ 0.75 V, below Vth(min) 1.0 V.
\item
  The ramp is then set by SDRV's 50--100 nA into C720: dVs/dt ≈ 5--10 V/s, so ≈ 10--15 mA into the bulk and VSYS reaches ≈ 7 V in about 1 s.
\item
  Bench gate for fitting (S2): 10 pack insertions with IBAT and SFET\_\allowbreak{}G scoped; peak \textless{} 6 A; no R5460 trip.
\end{itemize}

\paragraph{2.3 Design calculations}

\begin{itemize}
\tightlist
\item
  \textbf{PROG = 8.2 kΩ ±1 \%} selects 2S and 750 kHz at POR (Table 7-1: 1.5 MHz/2S = 6.04 k, 750 kHz/2S = 8.2 k).

  \begin{itemize}
  \tightlist
  \item
    \textbf{BQ25792 POR and watchdog defaults for 2S are VREG 8.4 V, ICHG 2 A, VSYSMIN 7.0 V} (BQ25792 ZHCSMD3C Table 9-2; the TI SLUSDG1C online edition agrees). The 1 A in BQ25798 Table 7-2 does not apply to the BQ25792 (release review S3).
  \item
    Blank, hung or watchdog-expired firmware therefore charges at 2 A: 0.4C on the 5 Ah pack, but \textbf{0.57C on a 3.5 Ah fallback pack} (over the 0.5C rule).
  \item
    \textbf{Firmware writes ICHG at every boot} (≤ 2.5 A for 5 Ah, ≤ 1.75 A for 3.5 Ah) and after every watchdog expiry (§7). If the 3.5 Ah pack is used, record the watchdog policy.
  \item
    Over-discharge recovery (R5460 DOUT low, cells 3.0--3.2 V): the 2 A default flows through Q801's body diode, ≈0.7 V × 2 A ≈ \textbf{1.5 W} for tens of minutes. Layout gives Q801 copper for 1.5 W (§6). Bench: recover a deep-discharged pack with a thermocouple on Q801.
  \end{itemize}
\item
  \textbf{Inductor (§8.2.2.2).} ISW(max) = 5 A (ROC §6.3). Ripple ΔI = VSYS·(1 − VSYS/VBUS)/(L·fsw) at VSYS = 8.4 V, L = 2.2 µH, fsw = 750 kHz: 15 V → 2.24 A, 20 V → 2.95 A, 12 V → 1.53 A. Boost mode at 5 V in → 1.23 A. Peak = 5 + 2.95/2 = \textbf{6.5 A} \textless{} Isat 10 A. Datasheet minimum is ``Isat ≥ 2 A for PFM pulses''. In practice the 15 V/2.5 A charge-while-play case peaks at ≈ 5 + 1.1 = 6.1 A.
\item
  \textbf{ILIM\_\allowbreak{}HIZ hardware clamp} (§7.3.4.3): V = 1 V + 0.8 Ω·ILIM. With REGN 4.8--5.0 V (EC table) and divider 100/(51+100) = 0.662: V = 3.18--3.31 V → \textbf{ILIM = 2.72--2.89 A} (3.05 A at REGN 5.2 V max). That keeps a 3 A PD contract within rating. The pin also sets the ceiling the host can program while EN\_\allowbreak{}EXTILIM = 1. Tie to REGN to remove it.
\item
  \textbf{D+/D− left open.} USB data belongs to the CH340C, and the BQ's BC1.2 source drive (0.6 V on D+) would disturb enumeration. With open D+/D−, the auto detection ends as SDP (500 mA) or ``unknown'' (3 A) (Table 7-4). Either is clamped by ILIM\_\allowbreak{}HIZ, and VINDPM (VBUS0 − 1.4 V ≥ 7 V source, §7.3.4.4) protects weak sources. \textbf{Firmware:} set AUTO\_\allowbreak{}INDET\_\allowbreak{}EN = 0, read CH224Q reg 0x50, and write IINDPM = min(PD current, 3.0 A). Without PD (5 V legacy), use ICO (EN\_\allowbreak{}ICO = 1) or 1.5 A.
\item
  \textbf{TS network} (Fig 7-12, §7.3.9.5). The thermistor is a 10 k NTC, B25/50 = 3435 (= 103AT). RT1/RT2 are E24 Basic values close to TI's 5.24 k/30.31 k (EVM: 5.23 k/30.1 k). Thresholds are VT1 = 73.3 \% REGN (cold stop) and VT5 = 34.2 \% (hot stop).

  \begin{itemize}
  \tightlist
  \item
    Computed with a β = 3435 model: \textbf{cold stop at 2.0 °C, hot stop at 60.8 °C} (TI values give 1.0/59.9 °C).
  \item
    With a B3380 NTC: 1.6/61.5 °C.
  \item
    Open NTC → TS = 30/35.1 = 85.5 \% \textgreater{} 73.3 \%, which reads as ``cold'' and suspends charging (fail-safe).
  \item
    Firmware sets JEITA T2/T3 (TS\_\allowbreak{}COOL/TS\_\allowbreak{}WARM) as desired.
  \end{itemize}
\item
  \textbf{Capacitors.} VBUS/PMID parts are 50 V (≥ 20 V PD headroom, TI asks for ≥ 25 V); effective ≈ 3.5--4 µF each at 15--20 V (TI SLUSEA4D quotes 3.45 µF for a 10 µF/50 V 1206 at 24 V), so PMID = 3 × ≈4 = ≈12 µF ≥ 4 µF min (ROC). SYS/BAT parts are 25 V for a ≤ 9 V rail (≥ 16 V required for 2S): ≈ 6 µF each effective, so SYS ≈ 30 µF ≥ 6 µF min, BAT ≈ 12 µF ≥ 3 µF min.
\item
  \textbf{Ship FET (Q701).} SDRV is a charge pump that drives the gate 5 V above BAT (§7.3.12), giving Vgs ≈ 5 V → Rds ≈ 2.8 mΩ. BATP is on the battery side (via R706), so the IC still sees the pack in ship mode. Modes (SDRV\_\allowbreak{}CTRL):

  \begin{itemize}
  \tightlist
  \item
    00 = on (POR default).
  \item
    01 = shutdown: I2C off, wake only by adapter.
  \item
    10 = ship: I2C on, wake by QON low ≥ tSM\_\allowbreak{}EXIT (1 s or 15 ms), adapter, or register write.
  \item
    11 = 350 ms system power reset.
  \end{itemize}

  Ship/shutdown are only accepted with no adapter. \textbf{Firmware must set SFET\_\allowbreak{}PRESENT = 1 and EN\_\allowbreak{}BATOC = 1} to enable the 9.3 A battery over-current cut-off (§7.3.13.1.7). \textbf{EN\_\allowbreak{}BATOC is reset by a watchdog expiry and by REG\_\allowbreak{}RST} (BQ25798 Table 7-30), so the full contract is in §7 (S1). Battery-side drain in ship mode is ≈ 15 µA: R5460 4 µA, 2 × HY2213 2.5 µA, mid divider 4.2 µA, BQ ship-mode quiescent. A 5 Ah pack lasts years.
\item
  \textbf{QON} (SW701 on the main board, or the UI-board button through the FFC): press ≥ 1 s (or 15 ms with WKUP\_\allowbreak{}DLY) to exit ship mode; hold ≥ 10 s for a full system power reset (§7.3.12.3). The 10 s reset is also the only battery-side reset for a hung ESP32 with the enclosure closed (S4).
\item
  \textbf{/CE} tied low via 10 k: charging starts without firmware (needed for dead-battery recovery). Firmware gates charging with EN\_\allowbreak{}CHG.
\item
  \textbf{STAT LED:} I = (4.8 − 2.8) V / 1 k ≈ 2 mA from REGN. REGN only exists with VBUS (or with the ADC on), so the LED never drains the battery. STAT sinks ≤ 6 mA (abs max). Low = charging; blinks 1 Hz on a fault.
\item
  \textbf{Battery current limit reminder.} BATFET is 6 A RMS continuous, 10 A ≤ 1 s, 11 mΩ (ROC §6.3). The boost ILIM (§4.1) plus the amp DRC must keep the average battery current ≤ 6 A.
\end{itemize}

\begin{center}\rule{0.5\linewidth}{0.5pt}\end{center}

\subsubsection{\texorpdfstring{3. Sheet \texttt{battery} (pack connectors, fuse, 2S protection, balancing, NTC)}{3. Sheet battery (pack connectors, fuse, 2S protection, balancing, NTC)}}

\paragraph{3.1 Component table}

{\def\LTcaptype{none} % do not increment counter
\begin{longtable}[]{@{}
  >{\raggedright\arraybackslash}p{(\linewidth - 14\tabcolsep) * \real{0.1250}}
  >{\raggedright\arraybackslash}p{(\linewidth - 14\tabcolsep) * \real{0.1250}}
  >{\raggedright\arraybackslash}p{(\linewidth - 14\tabcolsep) * \real{0.1250}}
  >{\raggedright\arraybackslash}p{(\linewidth - 14\tabcolsep) * \real{0.1250}}
  >{\raggedright\arraybackslash}p{(\linewidth - 14\tabcolsep) * \real{0.1250}}
  >{\raggedright\arraybackslash}p{(\linewidth - 14\tabcolsep) * \real{0.1250}}
  >{\raggedright\arraybackslash}p{(\linewidth - 14\tabcolsep) * \real{0.1250}}
  >{\raggedright\arraybackslash}p{(\linewidth - 14\tabcolsep) * \real{0.1250}}@{}}
\toprule\noalign{}
\begin{minipage}[b]{\linewidth}\raggedright
Ref
\end{minipage} & \begin{minipage}[b]{\linewidth}\raggedright
Value
\end{minipage} & \begin{minipage}[b]{\linewidth}\raggedright
Footprint
\end{minipage} & \begin{minipage}[b]{\linewidth}\raggedright
Rating
\end{minipage} & \begin{minipage}[b]{\linewidth}\raggedright
LCSC
\end{minipage} & \begin{minipage}[b]{\linewidth}\raggedright
B/E
\end{minipage} & \begin{minipage}[b]{\linewidth}\raggedright
Symbol
\end{minipage} & \begin{minipage}[b]{\linewidth}\raggedright
Notes
\end{minipage} \\
\midrule\noalign{}
\endhead
\bottomrule\noalign{}
\endlastfoot
J801 & XT30PW-M (male, PCB R/A) & \texttt{Connector\_\allowbreak{}AMASS:\allowbreak{}AMASS\_\allowbreak{}XT30PW-M\_\allowbreak{}1x02\_\allowbreak{}P2.50mm\_\allowbreak{}Horizontal} & 15 A, 1.2 mΩ & C431092 & E & \texttt{Connector\_\allowbreak{}Generic:\allowbreak{}Conn\_\allowbreak{}01x02} & \textbf{Pin 2 = +, pin 1 = −} (the footprint silkscreen shows ``+'' next to pad 2 at x = −5 and ``−'' next to pad 1 at x = 0) \\
J802 & B3B-XH-A (JST-XH 3P vertical) & \texttt{Connector\_\allowbreak{}JST:\allowbreak{}JST\_\allowbreak{}XH\_\allowbreak{}B3B-XH-A\_\allowbreak{}1x03\_\allowbreak{}P2.50mm\_\allowbreak{}Vertical} & 3 A & C144394 & E & \texttt{Connector\_\allowbreak{}Generic:\allowbreak{}Conn\_\allowbreak{}01x03} & Balance lead \\
J803 & B2B-PH-K-S (JST-PH 2P) & \texttt{Connector\_\allowbreak{}JST:\allowbreak{}JST\_\allowbreak{}PH\_\allowbreak{}B2B-PH-K\_\allowbreak{}1x02\_\allowbreak{}P2.00mm\_\allowbreak{}Vertical} & 2 A & C131337 & E & \texttt{Connector\_\allowbreak{}Generic:\allowbreak{}Conn\_\allowbreak{}01x02} & Pack NTC (\textbf{default sensor}, S6). Place at the pack-lead exit \\
F801 & 10 A very fast, 0451010.MRL & \textbf{custom} \texttt{power:\allowbreak{}Fuse\_\allowbreak{}2410\_\allowbreak{}6125Metric} (§0) & 125 V, 5.6 mΩ, I²t 26.5 A²s; derate to 7.5 A continuous & C44479 & E & \texttt{Device:\allowbreak{}Fuse} & Pack + \\
U801 & R5460N212AF & SOT-23-6 & OV 4.29 V (rel 4.05), UV 3.00 V (rel 3.20), excess-discharge 0.200 V/12 ms, short 1.1 V/300 µs & C259717 & E & custom \texttt{power:\allowbreak{}R5460x2} & 2S protector \\
R801 & 330 Ω & R\_\allowbreak{}0603 & & C23138 & B & \texttt{Device:\allowbreak{}R} & VDD supply (Ricoh R1) \\
C801 & 100 nF & C\_\allowbreak{}0603 & 50 V & C14663 & B & \texttt{Device:\allowbreak{}C} & VDD--VSS (C1) \\
R802 & 330 Ω & R\_\allowbreak{}0603 & & C23138 & B & \texttt{Device:\allowbreak{}R} & VC (R2) \\
C802 & 100 nF & C\_\allowbreak{}0603 & 50 V & C14663 & B & \texttt{Device:\allowbreak{}C} & VC--VSS (C2) \\
C803 & 10 nF & C\_\allowbreak{}0603 & 50 V & C57112 & B & \texttt{Device:\allowbreak{}C} & VSS--V− (C3, ``must mount ≥ 0.01 µF'') \\
R803 & 1.0 kΩ & R\_\allowbreak{}0603 & & C21190 & B & \texttt{Device:\allowbreak{}R} & V− sense (R3; R1 + R3 ≥ 1 k, R3 ≤ 3 k) \\
Q801 & NCEP3065QU & \texttt{Package\_\allowbreak{}SON:\allowbreak{}VSON-8\_\allowbreak{}3.3x3.3mm\_\allowbreak{}P0.65mm\_\allowbreak{}NexFET} & 30 V 65 A, 2.5 mΩ@6 V est. & C502964 & E & \texttt{Transistor\_\allowbreak{}FET:\allowbreak{}CSD17581Q3A} & Discharge FET (DOUT) \\
Q802 & NCEP3065QU & same & same & C502964 & E & same & Charge FET (COUT) \\
U802 & HY2213-BB3A & SOT-23-6 & VCU 4.200 ± 0.025 V, VCR 4.190 V, N-FET drive & C113632 & E & custom \texttt{power:\allowbreak{}HY2213} & Lower cell balancer \\
U803 & HY2213-BB3A & SOT-23-6 & & C113632 & E & custom \texttt{power:\allowbreak{}HY2213} & Upper cell balancer \\
R804, R807 & 100 Ω & R\_\allowbreak{}0603 & & C22775 & B & \texttt{Device:\allowbreak{}R} & HY2213 VDD (HY2213 §9: 100--200 Ω) \\
C804, C805 & 100 nF & C\_\allowbreak{}0603 & 50 V & C14663 & B & \texttt{Device:\allowbreak{}C} & HY2213 VDD--VSS \\
Q803, Q804 & AO3400A & \texttt{Package\_\allowbreak{}TO\_\allowbreak{}SOT\_\allowbreak{}SMD:\allowbreak{}SOT-23} & 30 V, Rds 48 mΩ@2.5 V & C20917 & B & \texttt{Transistor\_\allowbreak{}FET:\allowbreak{}AO3400A} & Balance switches \\
R805, R806, R808, R809 & 120 Ω 1 \% & \texttt{Resistor\_\allowbreak{}SMD:\allowbreak{}R\_\allowbreak{}1206\_\allowbreak{}3216Metric} & \textbf{0.25 W} each & C17909 & B & \texttt{Device:\allowbreak{}R} & Bleed: 2 in parallel per cell \\
R810, R811 & 1.0 MΩ & R\_\allowbreak{}0603 & & C22935 & B & \texttt{Device:\allowbreak{}R} & Mid-tap bias when the balance lead is unplugged \\
TH801 & 10 kΩ NTC, B25/50 = 3435 K, 1 \% & \texttt{Resistor\_\allowbreak{}SMD:\allowbreak{}R\_\allowbreak{}0603\_\allowbreak{}1608Metric} & & C2889048 & E & \texttt{Device:\allowbreak{}Thermistor\_\allowbreak{}NTC} & On-board NTC (CMFA103F3435). \textbf{Alternate} sensor; fitted but unused in the default jumper position \\
JP801 & open & \texttt{Jumper:\allowbreak{}SolderJumper-2\_\allowbreak{}P1.3mm\_\allowbreak{}Open\_\allowbreak{}RoundedPad1.0x1.5mm} & & -- & -- & \texttt{Jumper:\allowbreak{}SolderJumper\_\allowbreak{}2\_\allowbreak{}Open} & balance B− Kelvin (default open) \\
JP802 & open & same & & -- & -- & \texttt{Jumper:\allowbreak{}SolderJumper\_\allowbreak{}2\_\allowbreak{}Open} & balance B+ Kelvin (default open) \\
JP803 & bridged 1--2 & \texttt{Jumper:\allowbreak{}SolderJumper-3\_\allowbreak{}P1.3mm\_\allowbreak{}Bridged12\_\allowbreak{}RoundedPad1.0x1.5mm} & & -- & -- & \texttt{Jumper:\allowbreak{}SolderJumper\_\allowbreak{}3\_\allowbreak{}Bridged12} & NTC select (ECO-002 \#3): \textbf{1--2 = pack NTC on J803 (default)}; cut and bridge 2--3 for on-board TH801. KiCad 10 has no stock Bridged23 jumper, so the outer-pad nets are swapped instead \\
TP801 & VBAT & \texttt{TestPoint:\allowbreak{}TestPoint\_\allowbreak{}Pad\_\allowbreak{}D1.5mm} & & -- & -- & \texttt{Connector:\allowbreak{}TestPoint} & \texttt{VBAT\_\allowbreak{}PACK} \\
TP802 & BATT\_\allowbreak{}N & same & & -- & -- & \texttt{Connector:\allowbreak{}TestPoint} & cell negative \\
TP803 & BAL\_\allowbreak{}MID & same & & -- & -- & \texttt{Connector:\allowbreak{}TestPoint} & \\
\end{longtable}
}

\paragraph{3.2 Pin-by-pin}

\textbf{J801 XT30PW-M}: pin 2 (+) \texttt{PACK\_\allowbreak{}P}; pin 1 (−) \texttt{BATT\_\allowbreak{}N}. The two unnumbered mounting holes are mechanical (no net).

\textbf{J802 JST-XH 2S balance} (assumes the standard 2S lead: black 0 V, middle tap, red 8.4 V; \textbf{verify the pack's wire order}): pin 1 \texttt{BAL\_\allowbreak{}N} → JP801 pin 1; JP801 pin 2 \texttt{BATT\_\allowbreak{}N}. Pin 2 \texttt{BAL\_\allowbreak{}MID}. Pin 3 \texttt{BAL\_\allowbreak{}P} → JP802 pin 1; JP802 pin 2 \texttt{PACK\_\allowbreak{}P} (before the fuse, so the fuse still protects).

\textbf{F801}: pin 1 \texttt{PACK\_\allowbreak{}P}, pin 2 \texttt{VBAT\_\allowbreak{}PACK} (global).

\textbf{U801 R5460N212AF} (Ricoh typical application, p.18):

{\def\LTcaptype{none} % do not increment counter
\begin{longtable}[]{@{}
  >{\raggedright\arraybackslash}p{(\linewidth - 6\tabcolsep) * \real{0.2500}}
  >{\raggedright\arraybackslash}p{(\linewidth - 6\tabcolsep) * \real{0.2500}}
  >{\raggedright\arraybackslash}p{(\linewidth - 6\tabcolsep) * \real{0.2500}}
  >{\raggedright\arraybackslash}p{(\linewidth - 6\tabcolsep) * \real{0.2500}}@{}}
\toprule\noalign{}
\begin{minipage}[b]{\linewidth}\raggedright
Pin
\end{minipage} & \begin{minipage}[b]{\linewidth}\raggedright
Name
\end{minipage} & \begin{minipage}[b]{\linewidth}\raggedright
Net
\end{minipage} & \begin{minipage}[b]{\linewidth}\raggedright
Connection
\end{minipage} \\
\midrule\noalign{}
\endhead
\bottomrule\noalign{}
\endlastfoot
1 & DOUT & \texttt{PROT\_\allowbreak{}DO} & Q801 gate \\
2 & COUT & \texttt{PROT\_\allowbreak{}CO} & Q802 gate \\
3 & V− & \texttt{PROT\_\allowbreak{}VM} & R803 1 k to \texttt{GND} (= pack P−); C803 10 nF to \texttt{BATT\_\allowbreak{}N} \\
4 & VC & \texttt{PROT\_\allowbreak{}VC} & R802 330 Ω from \texttt{BAL\_\allowbreak{}MID}; C802 100 nF to \texttt{BATT\_\allowbreak{}N} \\
5 & VDD & \texttt{PROT\_\allowbreak{}VDD} & R801 330 Ω from \texttt{VBAT\_\allowbreak{}PACK}; C801 100 nF to \texttt{BATT\_\allowbreak{}N} \\
6 & VSS & \texttt{BATT\_\allowbreak{}N} & cell negative \\
\end{longtable}
}

\textbf{Q801 (discharge FET)}: S (1,2,3) \texttt{BATT\_\allowbreak{}N}, G (4) \texttt{PROT\_\allowbreak{}DO}, D (5) \texttt{PROT\_\allowbreak{}D}.
\textbf{Q802 (charge FET)}: S (1,2,3) \texttt{GND}, G (4) \texttt{PROT\_\allowbreak{}CO}, D (5) \texttt{PROT\_\allowbreak{}D}. The drains are common (back-to-back), as in the Ricoh figure.

\textbf{U802 HY2213 (lower cell, between BATT\_\allowbreak{}N and BAL\_\allowbreak{}MID)}: 1 NC, 2 VDD \texttt{BAL1\_\allowbreak{}VDD} (R804 100 Ω from \texttt{BAL\_\allowbreak{}MID}; C804 to \texttt{BATT\_\allowbreak{}N}), 3 VSS \texttt{BATT\_\allowbreak{}N}, 4 NC, 5 NC, 6 OUT \texttt{BAL1\_\allowbreak{}G}.
\textbf{Q803}: G \texttt{BAL1\_\allowbreak{}G}, S \texttt{BATT\_\allowbreak{}N}, D \texttt{BAL1\_\allowbreak{}D}. \textbf{R805 ∥ R806} (120 Ω each): \texttt{BAL\_\allowbreak{}MID} ↔ \texttt{BAL1\_\allowbreak{}D}.

\textbf{U803 HY2213 (upper cell, between BAL\_\allowbreak{}MID and VBAT\_\allowbreak{}PACK)}: 2 VDD \texttt{BAL2\_\allowbreak{}VDD} (R807 100 Ω from \texttt{VBAT\_\allowbreak{}PACK}; C805 to \texttt{BAL\_\allowbreak{}MID}), 3 VSS \texttt{BAL\_\allowbreak{}MID}, 6 OUT \texttt{BAL2\_\allowbreak{}G}, 1/4/5 NC.
\textbf{Q804}: G \texttt{BAL2\_\allowbreak{}G}, S \texttt{BAL\_\allowbreak{}MID}, D \texttt{BAL2\_\allowbreak{}D}. \textbf{R808 ∥ R809}: \texttt{VBAT\_\allowbreak{}PACK} ↔ \texttt{BAL2\_\allowbreak{}D}.

\textbf{R810}: \texttt{VBAT\_\allowbreak{}PACK} ↔ \texttt{BAL\_\allowbreak{}MID}. \textbf{R811}: \texttt{BAL\_\allowbreak{}MID} ↔ \texttt{BATT\_\allowbreak{}N}.

\textbf{NTC selection}: TH801 pin 1 \texttt{BAT\_\allowbreak{}NTC\_\allowbreak{}INT}, pin 2 \texttt{GND}. J803 pin 1 \texttt{BAT\_\allowbreak{}NTC\_\allowbreak{}EXT}, pin 2 \texttt{GND}. JP803: \textbf{pin 1 (A) \texttt{BAT\_\allowbreak{}NTC\_\allowbreak{}EXT}}, pin 2 (C, common) \texttt{BAT\_\allowbreak{}TS} (global), \textbf{pin 3 (B) \texttt{BAT\_\allowbreak{}NTC\_\allowbreak{}INT}}. With the factory bridge 1--2, the pack NTC is used.

\paragraph{3.3 Full current path (pack → BQ BAT pin) and calculations}

{\scriptsize
\begin{verbatim}
Pack+ ─XT30 pin2─ PACK_P ─F801 10A─ VBAT_PACK ─Q701 ship FET (D→S)─ CHG_BAT ─ U701 BAT(22,23) ─BATFET 11mΩ─ SYS = VSYS
Pack− ─XT30 pin1─ BATT_N ─Q801 (S→D)─ PROT_D ─Q802 (D→S)─ GND (= every ground on the board)
\end{verbatim}
}

Series resistance ≈ connector 2 × 1.2 + fuse 5.6 + Q701 2.8 + 2 × 2.5 (Q801/Q802 at Vgs ≈ VDD ≈ 7 V) + BATFET 11 ≈ \textbf{27 mΩ} → 0.19 V and 1.3 W total at 7 A, 0.97 W at 6 A. Q801/Q802 each dissipate 2.5 mΩ × 6² = 90 mW.

\begin{itemize}
\tightlist
\item
  \textbf{Protection thresholds (R5460N212AF, selection table).} Over-charge 4.290 V (rel 4.050 V, 1.0 s); over-discharge 3.000 V (rel 3.200 V, 128 ms). ``AF'' version: over-discharge releases only after a charger is connected, then at the release voltage. Standby ≤ 0.1 µA.

  \begin{itemize}
  \tightlist
  \item
    Excess discharge current: V− \textgreater{} 0.200 V for 12 ms → 0.2 V / (2 × 2.5 mΩ) ≈ \textbf{40 A}.
  \item
    Short: V− \textgreater{} 1.1 V for 300 µs.
  \end{itemize}

  Both are last-resort only. Real over-current protection comes from the BQ IBAT\_\allowbreak{}OCP 9.3 A (with Q701), the boost ILIM (§4.1) and F801.
\item
  \textbf{Gate drive.} DOUT/COUT are CMOS outputs swinging VDD (6.0--8.4 V) → Vgs ≈ 6--8 V on the NCEP3065QU (Vth 1.0--2.0 V, Rds 3.0 mΩ@4.5 V, 1.9 mΩ@10 V). No gate resistors (Ricoh figure).
\item
  \textbf{Fuse.} 10 A rating, derated 25 \% → 7.5 A continuous at 25 °C (Littelfuse note), which is ≥ the 6 A design ceiling. 200 \% opens ≤ 5 s.
\item
  \textbf{Balancing (HY2213 §6, §9).} HY2213-BB3A switches on above 4.200 V and off below 4.190 V (code A = N-MOSFET drive, OUT active high). Bleed R = 120 ∥ 120 = 60 Ω → I = 4.2/60 = \textbf{70 mA}. Each 1206 dissipates 4.2²/120 = \textbf{0.147 W} (59 \% of 0.25 W). That is ≈ 0.3 W per balancing cell, only at the very end of a charge. HY2213 notes: R(VDD) 100--200 Ω, C 0.01--1 µF.
\item
  \textbf{Mid-tap bias.} R810/R811 hold \texttt{BAL\_\allowbreak{}MID} = VBAT/2 when the balance plug is out, so neither the R5460 nor the HY2213s see a floating cell (drain 8.4 V / 2 MΩ = 4.2 µA). The real cell junction overrides it when J802 is plugged.

  \begin{itemize}
  \tightlist
  \item
    \textbf{Residual risk (S7, accepted):} with J802 unplugged, \textbf{cell-level OV/UV protection is silently lost}. For example, 4.45 V + 3.95 V reads as 4.2 V + 4.2 V. There is no ADC on BAL\_\allowbreak{}MID to detect this.
  \item
    \textbf{Assembly step, mandatory before the first charge and after any pack swap:}

    \begin{enumerate}
    \def\labelenumi{\arabic{enumi}.}
    \tightlist
    \item
      Check J802 is fully seated.
    \item
      With the pack connected, measure TP803 − TP802. It should read the lower cell voltage, 3.0--4.2 V.
    \item
      Briefly load TP803--TP802 with a 1 kΩ resistor. A seated lead holds the voltage (drop \textless{} 10 mV). An unplugged lead collapses toward 0 V, because the 1 MΩ divider can't source current.
    \end{enumerate}

    Record the result in the build log.
  \end{itemize}
\item
  \textbf{NTC (S6).} Default: a 10 kΩ, B ≈ 3435, \textbf{isolated 2-wire} NTC bonded to the pack face and wired to J803 (the pack itself has no NTC). Record it in the pack harness drawing.

  \begin{itemize}
  \tightlist
  \item
    The pack sits on a tray 7 mm above the board, so TH801, next to J801, would track board temperature, not cell temperature. TH801 remains as the alternate.
  \item
    Don't use a pack NTC that is referenced to the cell negative inside the pack, because B− ≠ GND here.
  \item
    With J803 unplugged, TS reads ``cold'' (85.5 \% of REGN) and \textbf{charging is suspended}. This is fail-safe, but the harness NTC must be fitted for the pack to charge.
  \item
    The TS window covers \textbf{charging only}. There is no hardware discharge-temperature cut-off, so firmware must read TS in battery mode and stop playback above the pack limit (§7).
  \end{itemize}
\end{itemize}

\begin{center}\rule{0.5\linewidth}{0.5pt}\end{center}

\subsubsection{\texorpdfstring{4. Sheet \texttt{rails}}{4. Sheet rails}}

\paragraph{4.1 TPS61178 boost, VSYS (6.0--8.6 V) → PVDD 18.0 V}

\subparagraph{4.1.1 Component table}

{\def\LTcaptype{none} % do not increment counter
\begin{longtable}[]{@{}
  >{\raggedright\arraybackslash}p{(\linewidth - 14\tabcolsep) * \real{0.1250}}
  >{\raggedright\arraybackslash}p{(\linewidth - 14\tabcolsep) * \real{0.1250}}
  >{\raggedright\arraybackslash}p{(\linewidth - 14\tabcolsep) * \real{0.1250}}
  >{\raggedright\arraybackslash}p{(\linewidth - 14\tabcolsep) * \real{0.1250}}
  >{\raggedright\arraybackslash}p{(\linewidth - 14\tabcolsep) * \real{0.1250}}
  >{\raggedright\arraybackslash}p{(\linewidth - 14\tabcolsep) * \real{0.1250}}
  >{\raggedright\arraybackslash}p{(\linewidth - 14\tabcolsep) * \real{0.1250}}
  >{\raggedright\arraybackslash}p{(\linewidth - 14\tabcolsep) * \real{0.1250}}@{}}
\toprule\noalign{}
\begin{minipage}[b]{\linewidth}\raggedright
Ref
\end{minipage} & \begin{minipage}[b]{\linewidth}\raggedright
Value
\end{minipage} & \begin{minipage}[b]{\linewidth}\raggedright
Footprint
\end{minipage} & \begin{minipage}[b]{\linewidth}\raggedright
Rating
\end{minipage} & \begin{minipage}[b]{\linewidth}\raggedright
LCSC
\end{minipage} & \begin{minipage}[b]{\linewidth}\raggedright
B/E
\end{minipage} & \begin{minipage}[b]{\linewidth}\raggedright
Symbol
\end{minipage} & \begin{minipage}[b]{\linewidth}\raggedright
Notes / citation
\end{minipage} \\
\midrule\noalign{}
\endhead
\bottomrule\noalign{}
\endlastfoot
U901 & TPS61178RNWR & \textbf{custom} \texttt{power:\allowbreak{}Texas\_\allowbreak{}RNW0013A\_\allowbreak{}VQFN-HR-13\_\allowbreak{}3x3.5mm\_\allowbreak{}P0.5mm} & Vin 2.7--20 V, Vout 4.5--20 V, OVP 21 V (20.5 min) & C127918 & E & \textbf{custom} \texttt{power:\allowbreak{}TPS61178} & \\
L901 & 2.2 µH FXL1040-2R2-M & \texttt{Inductor\_\allowbreak{}SMD:\allowbreak{}L\_\allowbreak{}Changjiang\_\allowbreak{}FXL1040} & Irms 12 A, Isat 18 A, 7 mΩ & C167228 & E & \texttt{Device:\allowbreak{}L} & \\
C901, C902 & 22 µF & C\_\allowbreak{}1206 & 25 V X5R & C12891 & B & \texttt{Device:\allowbreak{}C} & VIN bulk (VSYS ≤ 9 V) \\
C903 & 100 nF & C\_\allowbreak{}0603 & 50 V & C14663 & B & \texttt{Device:\allowbreak{}C} & VIN pin 12 \\
C904 & 4.7 µF & C\_\allowbreak{}0603 & 16 V X5R & C19666 & B & \texttt{Device:\allowbreak{}C} & VCC (≥ 10 × CBST, §9.2.4.4.5) \\
C905 & 100 nF & C\_\allowbreak{}0603 & 50 V X7R & C14663 & B & \texttt{Device:\allowbreak{}C} & BST--SW (0.1--1 µF, ≥ 10 V, §9.2.4.4.4) \\
C906--C913 & 10 µF (×8) & C\_\allowbreak{}1206 & \textbf{50 V} X5R & C13585 & B & \texttt{Device:\allowbreak{}C} & COUT on \texttt{PVDD} at pins 8/6 (≈ 4.5 µF each at 18 V → ≈ 36 µF) \\
C914 & 100 nF & C\_\allowbreak{}0603 & 50 V & C14663 & B & \texttt{Device:\allowbreak{}C} & COUT1 HF, at pins 8/6 \\
R901 & 300 kΩ 1 \% & R\_\allowbreak{}0603 & & C23024 & B & \texttt{Device:\allowbreak{}R} & FREQ → AGND: 599 kHz \\
R902 & 82 kΩ 1 \% & R\_\allowbreak{}0603 & & C23254 & B & \texttt{Device:\allowbreak{}R} & ILIMIT → AGND \\
R903 & 100 kΩ 1 \% & R\_\allowbreak{}0603 & & C25803 & B & \texttt{Device:\allowbreak{}R} & Rc (ECO-001) \\
C915 & 15 nF & C\_\allowbreak{}0603 & 50 V X7R & C1596 & E & \texttt{Device:\allowbreak{}C} & Cc (ECO-001) \\
C916 & 10 pF & C\_\allowbreak{}0603 & 50 V C0G & C1634 & B & \texttt{Device:\allowbreak{}C} & Cp, COMP → AGND \\
R904 & 300 kΩ 1 \% & R\_\allowbreak{}0603 & & C23024 & B & \texttt{Device:\allowbreak{}R} & RUP: PVDD → FB \\
R905 & 24 kΩ 1 \% & R\_\allowbreak{}0603 & & C23352 & B & \texttt{Device:\allowbreak{}R} & RDOWN1: FB → AGND \\
R906 & 200 kΩ 1 \% & R\_\allowbreak{}0603 & & C25811 & B & \texttt{Device:\allowbreak{}R} & RDOWN2: FB → AGND \\
R907 & 68 kΩ 1 \% & R\_\allowbreak{}0603 & & C23231 & B & \texttt{Device:\allowbreak{}R} & trim injection, \textbf{fitted} (ECO-001) \\
D901 & 1N4148W & \texttt{Diode\_\allowbreak{}SMD:\allowbreak{}D\_\allowbreak{}SOD-123} & 75 V & C81598 & B & \texttt{Device:\allowbreak{}D} (1 K, 2 A) & trim diode, \textbf{fitted} (ECO-001) \\
R908, R909 & 10 kΩ & R\_\allowbreak{}0603 & & C25804 & B & \texttt{Device:\allowbreak{}R} & trim RC (fitted) \\
C917, C918 & 1 µF & C\_\allowbreak{}0603 & 50 V X5R & C15849 & B & \texttt{Device:\allowbreak{}C} & trim RC (fitted) \\
C920 & 470 µF & \texttt{Capacitor\_\allowbreak{}SMD:\allowbreak{}CP\_\allowbreak{}Elec\_\allowbreak{}10x10} & \textbf{35 V}, 105 °C, 10×10.2 mm (KNSCHA RVT470UF35V67RV0056) & C2836436 & E & \texttt{Device:\allowbreak{}C\_\allowbreak{}Polarized} & PVDD bulk (same part as the amp sheets) \\
TP901 & PVDD & \texttt{TestPoint:\allowbreak{}TestPoint\_\allowbreak{}Pad\_\allowbreak{}D2.0mm} & & -- & -- & \texttt{Connector:\allowbreak{}TestPoint} & \\
\end{longtable}
}

Pin types for the custom symbol are in §0.1.

\subparagraph{\texorpdfstring{4.1.2 Custom footprint \texttt{Texas\_\allowbreak{}RNW0013A\_\allowbreak{}VQFN-HR-13\_\allowbreak{}3x3.5mm\_\allowbreak{}P0.5mm} (TPS61178 §13, RNW0013A land pattern; top view, KiCad coordinates, +y down)}{4.1.2 Custom footprint Texas\_\allowbreak{}RNW0013A\_\allowbreak{}VQFN-HR-13\_\allowbreak{}3x3.5mm\_\allowbreak{}P0.5mm (TPS61178 §13, RNW0013A land pattern; top view, KiCad coordinates, +y down)}}

\begin{itemize}
\tightlist
\item
  Body 3.50 (x) × 3.00 (y) mm, 1.0 mm max height. Pin 1 at top-left.
\item
  Pads 1--5: roundrect 0.70 (x) × 0.25 (y) mm, r = 0.05. Centres x = −1.60; y = −1.00 (pin 1), −0.50 (2), 0.00 (3), +0.50 (4), +1.00 (5).
\item
  Pads 13--9: same size, x = +1.60; y = −1.00 (pin 13), −0.50 (12), 0.00 (11), +0.50 (10), +1.00 (9).
\item
  Pads 6, 7, 8 (power bars): roundrect 0.25 (x) × 3.40 (y) mm, centred at y = 0. x = −0.55 (pin 6 PGND), 0.00 (pin 7 SW), +0.55 (pin 8 VOUT). Solder-mask-defined in TI's pattern (mask opening ≈ pad + 0.05 mm per side). The bars extend 0.2 mm past the body top and bottom.
\item
  Courtyard 4.30 × 3.90 mm. Paste 1:1 on pads 1--5/9--13; split the bars into 3 apertures each (TI stencil ≈ 90 \%).
\end{itemize}

\subparagraph{4.1.3 U901 pin-by-pin}

{\def\LTcaptype{none} % do not increment counter
\begin{longtable}[]{@{}
  >{\raggedright\arraybackslash}p{(\linewidth - 6\tabcolsep) * \real{0.2500}}
  >{\raggedright\arraybackslash}p{(\linewidth - 6\tabcolsep) * \real{0.2500}}
  >{\raggedright\arraybackslash}p{(\linewidth - 6\tabcolsep) * \real{0.2500}}
  >{\raggedright\arraybackslash}p{(\linewidth - 6\tabcolsep) * \real{0.2500}}@{}}
\toprule\noalign{}
\begin{minipage}[b]{\linewidth}\raggedright
Pin
\end{minipage} & \begin{minipage}[b]{\linewidth}\raggedright
Name
\end{minipage} & \begin{minipage}[b]{\linewidth}\raggedright
Net
\end{minipage} & \begin{minipage}[b]{\linewidth}\raggedright
Connection
\end{minipage} \\
\midrule\noalign{}
\endhead
\bottomrule\noalign{}
\endlastfoot
1 & FREQ/SYNC & \texttt{BST\_\allowbreak{}FREQ} & R901 300 k to AGND (must not float) \\
2 & AGND & \texttt{GND} & quiet ground; tie to PGND at one point (pin 6 pad) \\
3 & ILIMIT & \texttt{BST\_\allowbreak{}ILIM} & R902 82 k to AGND \\
4 & COMP & \texttt{BST\_\allowbreak{}COMP} & R903 100 k → \texttt{BST\_\allowbreak{}CZ} → C915 15 nF → AGND; C916 10 pF COMP → AGND \\
5 & FB & \texttt{BST\_\allowbreak{}FB} & R904 from \texttt{PVDD}; R905 and R906 to AGND; R907 (fitted) to \texttt{TRIM\_\allowbreak{}K} \\
6 & PGND & \texttt{GND} & power ground, thermal vias \\
7 & SW & \texttt{BST\_\allowbreak{}SW} & L901 pin 2; C905 to \texttt{BST\_\allowbreak{}BOOT} \\
8 & VOUT & \texttt{PVDD} (global) & C906--C914 at the pin; C920; R904 sense \\
9 & DISDRV & no connect (KiCad no-connect flag) & Load-disconnect not used (§4.1.4) \\
10 & VCC & \texttt{BST\_\allowbreak{}VCC} & C904 4.7 µF to GND \\
11 & EN & \texttt{BOOST\_\allowbreak{}EN} (global) & from ESP32 GPIO2. 100 k pull-down is R107 on the mcu sheet. Internal 800 k pull-down too \\
12 & VIN & \texttt{VSYS} & C903 at the pin, C901/C902 \\
13 & BST & \texttt{BST\_\allowbreak{}BOOT} & C905 to \texttt{BST\_\allowbreak{}SW} \\
\end{longtable}
}

\textbf{L901}: pin 1 \texttt{VSYS}, pin 2 \texttt{BST\_\allowbreak{}SW}.
\textbf{Trim network}: R908 \texttt{PVDD\_\allowbreak{}TRIM} ↔ \texttt{TRIM\_\allowbreak{}A}; C917 \texttt{TRIM\_\allowbreak{}A} ↔ GND; R909 \texttt{TRIM\_\allowbreak{}A} ↔ \texttt{TRIM\_\allowbreak{}RC}; C918 \texttt{TRIM\_\allowbreak{}RC} ↔ GND; D901 anode (pin 2) \texttt{TRIM\_\allowbreak{}RC}, cathode (pin 1) \texttt{TRIM\_\allowbreak{}K}; R907 \texttt{TRIM\_\allowbreak{}K} ↔ \texttt{BST\_\allowbreak{}FB}.
\textbf{C920 (470 µF)}: + \texttt{PVDD}, − \texttt{GND}.

\subparagraph{4.1.4 Calculations (TPS61178 SLVSDA7E)}

\begin{itemize}
\tightlist
\item
  \textbf{Output voltage} (§9.2.4, eq. 4; VREF = 1.180/1.198/1.210 V min/typ/max):

  \begin{itemize}
  \tightlist
  \item
    RDOWN = 24 k ∥ 200 k = 21.43 k. PVDD = 1.198 × (1 + 300 k/21.43 k) = \textbf{17.97 V} (17.70--18.15 V over the VREF tolerance, ±1 \% resistors add ±0.2 V).
  \item
    Margin to OVP 20.5 V min and to Vout(max) 20 V: ≥ 1.85 V.
  \item
    Divider current 60 µA ≫ 100 × 60 nA FB leakage. RDOWN ≤ 100 k, as TI recommends.
  \end{itemize}
\item
  \textbf{PVDD\_\allowbreak{}TRIM} (diode-OR into FB; R907 and D901 \textbf{fitted}, ECO-001):

  \begin{itemize}
  \item
    Hand model (fixed 0.45 V diode drop): PVDD = 17.97 − 300 k × (Vt − 0.45 − 1.198)/(68 k + 20 k).
  \item
    SPICE (sim S4, 25 °C), which supersedes the hand values (at 5--15 µA the diode drops ≈0.35 V and the knee is soft):

    {\def\LTcaptype{none} % do not increment counter
    \begin{longtable}[]{@{}lll@{}}
    \toprule\noalign{}
    Vt & Hand model & SPICE \\
    \midrule\noalign{}
    \endhead
    \bottomrule\noalign{}
    \endlastfoot
    1.6 V & 17.97 V & 17.58 V \\
    2.0 V & 16.8 V & 16.43 V \\
    2.5 V & 15.1 V & 14.84 V \\
    3.0 V & 13.4 V & 13.20 V \\
    3.3 V & 12.3 V & 12.20 V \\
    \end{longtable}
    }

    Slope ≈ −3.4 V/V; range \textbf{18.0 → 12.2 V over 1.5--3.3 V}.
  \item
    \textbf{Hi-Z or 0 V ⇒ nominal 18 V.} The trim can only lower PVDD, apart from D901 reverse leakage: +2 mV at 25 °C, +19 mV at 60 °C, +86 mV at 85 °C (OVP 20.5 V min, harmless). A GPIO stuck high gives ≈12 V, and the amps still work.
  \item
    Two-pole RC filter (10 k/1 µF ×2): fc = 15.9 Hz. A ≥ 20 kHz LEDC PWM leaves \textless{} 1 µV of ripple at FB (sim). \textbf{A trim step settles in 70--110 ms; allow ≥ 150 ms.} The DAC1 output also works (8-bit = 44 mV of PVDD per LSB; 10-bit LEDC ≈ 11 mV).
  \item
    Temperature drift ≈ −8.6 mV/°C on PVDD while trimming. Firmware should calibrate duty → PVDD once.
  \end{itemize}
\item
  \textbf{Switching frequency} (§8.3.7, eq. 2--3): RFREQ = (1/f − 50 ns)/(3 × 1.8 pF). With 300 k: f = 1/(1.62 µs + 0.05 µs) = \textbf{599 kHz}. That keeps clear of the amps' 768 kHz and 1.536 MHz: 2 × 599 = 1.198 MHz (338 kHz from 1.536 MHz), 768 − 599 = 169 kHz. Every beat with m + n ≤ 7 is ≥ 91 kHz at nominal. The oscillator tolerance is ±20 \% (EC table at 342 k: 400--600 kHz), so high-order beats can land anywhere. Their amplitude is negligible, and they disappear in PFM. External sync is documented only for the forced-PWM TPS611781 (§8.3.8); that part is a drop-in swap (C2862009) if PFM bursts turn out audible.
\item
  \textbf{Current limit} (§8.3.5, eq. 1 RLIMIT{[}kΩ{]} = 745/ILIMIT{[}A{]}; EC table at 80.6 k: 6.4/8.0/9.4 A min/typ/max). 82 k → 9.1 A (eq. 1); scaled EC ≈ \textbf{6.3/7.9/9.2 A} (min/typ/max).

  \begin{itemize}
  \tightlist
  \item
    Load demand: at 40 W out, 18 V: Vin 6.0 V → Iin 7.0 A, ripple 3.04 A, peak 8.5 A. Vin 7.4 V → peak 7.3 A.
  \item
    So a low battery can current-limit at ≈ 30 W. That is intended: it keeps the battery ≤ 6 A RMS and below the BQ's 9.3 A IBAT\_\allowbreak{}OCP.
  \item
    The short-circuit limit (20 A) triggers hiccup (§8.3.6).
  \end{itemize}
\item
  \textbf{Inductor} (§9.2.4.1). ΔI = Vin·D/(L·f): 3.04 A (6 V), 3.31 A (7.4 V), 3.40 A (8.4 V). All are below TI's 4 A slope-compensation limit. Recommended range 0.47--3.3 µH. Peak ≤ 9.2 A (ILIM max) ≪ Isat 18 A; Irms ≤ 7.1 A \textless{} 12 A.
\item
  \textbf{Soft start into ≈ 1.5 mF.} Internal: a 2.6 ms pre-charge to 1.1 × VIN at 500 kHz, then a 3.2 ms reference ramp (§8.3.3). There is no soft-start pin. The hand estimate follows; \textbf{sim S1 (TI switching model) supersedes it} (see the end of this item).

  \begin{itemize}
  \tightlist
  \item
    PVDD is pre-biased at VSYS − 0.5 V (≈ 7 V = 39 \% of 18 V). That is above the 30 \% ``output short'' hiccup threshold, and pre-biased start-up is supported (§8.3.11).
  \item
    Pre-charge: 1.5 mF × ≈1.2 V / 2.6 ms ≈ 0.7 A. Benign.
  \item
    Ramp, 7 → 18 V in 3.2 ms: would need 1.5 mF × 11 V / 3.2 ms ≈ 5.2 A output (≈ 15 A input), which exceeds ILIM. The converter therefore runs in cycle-by-cycle current limit: average inductor current ≈ ILIM(min) 6.3 A − ½ × 3 A ripple ≈ 4.8 A, giving Iout ≈ 4.8 × 7/18 × 0.95 ≈ 1.8 A. PVDD reaches 18 V in ≈ \textbf{9 ms} (worst case, at minimum ILIM).
  \item
    Neither hiccup trigger is reached: inductor current is held at ILIM ≤ 9.2 A ≪ the 20 A short limit, and VOUT never falls below 30 \%. Battery current during this ≈ 10 ms is ≈ 5--8 A, within the BATFET's 10 A ≤ 1 s peak.
  \item
    \textbf{Sim S1 result (ECO-001):}

    \begin{itemize}
    \tightlist
    \item
      At 6.0 V and ILIM(min) the converter runs sub-harmonic in current limit (average IL 3.3--3.9 A), so the ramp takes \textbf{11 ms}.
    \item
      No hiccup, but only 0.39 V of margin to the 30 \% short threshold.
    \item
      The error amplifier winds up during the current-limited ramp, so \textbf{PVDD overshoots to 19.4 V (typ), and to the ≈21 V OVP worst case}. It stays there until the amps load it. All parts tolerate this (TAS5805M 26.4 V; caps 35/50 V), but it is avoided with the firmware soft start below.
    \end{itemize}
  \item
    \textbf{Firmware PVDD soft start (R907/D901 fitted):}

    \begin{enumerate}
    \def\labelenumi{\arabic{enumi}.}
    \tightlist
    \item
      Drive \texttt{PVDD\_\allowbreak{}TRIM} to 3.3 V (DAC 255 or 100 \% duty) \textbf{≥ 150 ms before} BOOST\_\allowbreak{}EN, so the RC is charged.
    \item
      BOOST\_\allowbreak{}EN high. PVDD rises only to ≈12.2 V, which keeps the converter out of current limit and prevents windup.
    \item
      ≥ 20 ms later, release amp PDN.
    \item
      Ramp \texttt{PVDD\_\allowbreak{}TRIM} down to the target over ≥ 30--50 ms. PVDD settles ≥ 150 ms after the last trim step.
    \end{enumerate}
  \item
    \textbf{Sequencing (both cases):} BOOST\_\allowbreak{}EN high → wait \textbf{≥ 20 ms} (covers the 11 ms ramp) → release amp PDN. Pull PDN low before dropping BOOST\_\allowbreak{}EN.
  \item
    Bench check: scope PVDD and IL at turn-on with all three 470 µF fitted, with and without the trim soft start.
  \end{itemize}
\item
  \textbf{Disconnect FET: not used.} DISDRV (pin 9) is left open: ``Leave it floating if not using the load disconnect function'' (pin table).

  \begin{itemize}
  \tightlist
  \item
    Why: the PVDD bulk is ≈ 3 × 470 µF + 9 × 10 µF ≈ 1.5 mF (this sheet + both amp sheets), but COUT1 can only be ≈ 36 µF. TI's ``COUT2 ≤ 10 × COUT1'' rule (§9.2.4.2) and the ≤ 3 ms FET turn-on limit (§9.2.4.4.3) can't be met, so a disconnect FET would risk start-up hiccup.
  \item
    Consequences: with EN low, PVDD = VSYS − Vf(body diode) ≈ 5.5--8 V. Amp PDN is low, so the draw is µA. TPS61178 shutdown current ≤ 3.5 µA; SW reverse leakage ≤ 6.5 µA.
  \item
    Inrush at battery connect: the bulk charges through Q701 as a source follower, ramped by SDRV's 50--100 nA into Crss 15.5--30 pF (higher at low V\_\allowbreak{}DS). That gives 1.7--6.5 V/ms, so \textbf{2.5--9.7 A} (release review S2; the v1.0 figure of 3.8 A assumed 40 pF). It is bounded only by fitting C720 together with C721 (rule and numbers in §2.2), after the bench check.
  \item
    With VBUS present, the charger's own soft start limits VSYS inrush.
  \item
    A hard PVDD short (battery → L901 → high-side body diode → short) bypasses the boost's hiccup protection. It is cleared by the BQ IBAT\_\allowbreak{}OCP ship-FET trip \textbf{only while the §7 firmware contract holds}. Otherwise only the R5460 (≈40 A for 12 ms) and F801 (slow at 10--20 A) act, and a 10--30 A soft fault can exceed the BATFET's 10 A/1 s rating (S1). Bench: 20 mΩ PVDD short with blank and with armed firmware. Record the trip time and peak IBAT. If the blank-firmware result is unacceptable, record the remedy (a lower-VDET3 R5460 variant, a smaller fuse, or restoring a disconnect) before routing.
  \end{itemize}
\item
  \textbf{Output capacitance.} COUT = C906--C913 (8 × 10 µF/50 V ≈ 36 µF effective at 18 V) + C914, all at pins 8/6. PVDD bulk: C920 (470 µF/35 V) on this sheet plus the amp sheets' 2 × 470 µF and 8 × 10 µF. At 18 V the 35 V electrolytic runs at 51 \% of rating.
\item
  \textbf{Loop compensation} (§9.2.4.4, eq. 15--29; ECO-001). \textbf{Rc = R903 = 100 kΩ, Cc = C915 = 15 nF, Cp = C916 = 10 pF.}

  \begin{itemize}
  \tightlist
  \item
    The zero is unchanged: 1/(2π × 100 k × 15 nF) = 106 Hz.
  \item
    Sim S1/S7 (1.48 mF bank, ESR 0.05--0.34 Ω, Vin 6.0--8.4 V): \textbf{PM ≥ 79°, GM ≥ 15 dB}.
  \item
    PVDD closed-loop impedance meets the 0.14 Ω target. A 0.2 → 2 A step at 6 V dips \textbf{0.34 V}.
  \item
    History: v1.1's 33 k/47 nF was stable (PM ≥ 66°) but left a 0.46--0.57 Ω impedance peak at ≈150 Hz (≈1 V dip, ≈2 V bass sag).
  \item
    Rejected earlier: TI's 15 k/6.8 nF and 22 k/10 nF (28--36° PM with ≈1.4 mF of low-ESR bulk).
  \item
    The 7.4 V load step was not simulated (ngspice stalled). Verify on the bench with 0.5 → 2.5 A at 6.0 and 7.4 V.
  \end{itemize}
\end{itemize}

\paragraph{4.2 TPS62933 buck, VSYS → +3V3}

\subparagraph{4.2.1 Component table}

{\def\LTcaptype{none} % do not increment counter
\begin{longtable}[]{@{}
  >{\raggedright\arraybackslash}p{(\linewidth - 14\tabcolsep) * \real{0.1250}}
  >{\raggedright\arraybackslash}p{(\linewidth - 14\tabcolsep) * \real{0.1250}}
  >{\raggedright\arraybackslash}p{(\linewidth - 14\tabcolsep) * \real{0.1250}}
  >{\raggedright\arraybackslash}p{(\linewidth - 14\tabcolsep) * \real{0.1250}}
  >{\raggedright\arraybackslash}p{(\linewidth - 14\tabcolsep) * \real{0.1250}}
  >{\raggedright\arraybackslash}p{(\linewidth - 14\tabcolsep) * \real{0.1250}}
  >{\raggedright\arraybackslash}p{(\linewidth - 14\tabcolsep) * \real{0.1250}}
  >{\raggedright\arraybackslash}p{(\linewidth - 14\tabcolsep) * \real{0.1250}}@{}}
\toprule\noalign{}
\begin{minipage}[b]{\linewidth}\raggedright
Ref
\end{minipage} & \begin{minipage}[b]{\linewidth}\raggedright
Value
\end{minipage} & \begin{minipage}[b]{\linewidth}\raggedright
Footprint
\end{minipage} & \begin{minipage}[b]{\linewidth}\raggedright
Rating
\end{minipage} & \begin{minipage}[b]{\linewidth}\raggedright
LCSC
\end{minipage} & \begin{minipage}[b]{\linewidth}\raggedright
B/E
\end{minipage} & \begin{minipage}[b]{\linewidth}\raggedright
Symbol
\end{minipage} & \begin{minipage}[b]{\linewidth}\raggedright
Notes
\end{minipage} \\
\midrule\noalign{}
\endhead
\bottomrule\noalign{}
\endlastfoot
U951 & TPS62933DRLR & \texttt{Package\_\allowbreak{}TO\_\allowbreak{}SOT\_\allowbreak{}SMD:\allowbreak{}SOT-583-8} & 3.8--30 V, 3 A, 12 µA IQ (PFM) & C3200405 & E & \texttt{Regulator\_\allowbreak{}Switching:\allowbreak{}TPS62933} & stock sym+fp \\
L951 & 4.7 µH FXL0420-4R7-M & \texttt{Inductor\_\allowbreak{}SMD:\allowbreak{}L\_\allowbreak{}Changjiang\_\allowbreak{}FXL0420} & Irms 2.8 A / Isat 3 A, 105 mΩ & C167208 & E & \texttt{Device:\allowbreak{}L} & Table 10-2 (3.3 V/500 kHz → 4.7 µH) \\
C951, C952 & 10 µF & C\_\allowbreak{}0805 & 25 V X5R & C15850 & B & \texttt{Device:\allowbreak{}C} & VIN \\
C953 & 100 nF & C\_\allowbreak{}0603 & 50 V & C14663 & B & \texttt{Device:\allowbreak{}C} & VIN HF, at pins 3--4 \\
R951 & 300 kΩ 1 \% & R\_\allowbreak{}0603 & & C23024 & B & \texttt{Device:\allowbreak{}R} & EN UVLO top (VSYS → EN) \\
R952 & 82 kΩ 1 \% & R\_\allowbreak{}0603 & & C23254 & B & \texttt{Device:\allowbreak{}R} & EN UVLO bottom \\
C954 & 33 nF & C\_\allowbreak{}0603 & 50 V X7R & C21117 & B & \texttt{Device:\allowbreak{}C} & SS (≥ 6.8 nF; do not float) \\
R953 & 0 Ω & R\_\allowbreak{}0603 & & C21189 & B & \texttt{Device:\allowbreak{}R} & BST series (TI: \textless{} 10 Ω allowed) \\
C955 & 100 nF & C\_\allowbreak{}0603 & 50 V X7R & C14663 & B & \texttt{Device:\allowbreak{}C} & BST--SW (≥ 16 V) \\
R954 & 75 kΩ 1 \% & R\_\allowbreak{}0603 & & C23242 & B & \texttt{Device:\allowbreak{}R} & FB top \\
R955 & 24 kΩ 1 \% & R\_\allowbreak{}0603 & & C23352 & B & \texttt{Device:\allowbreak{}R} & FB bottom \\
C956 & 10 pF & C\_\allowbreak{}0603 & 50 V C0G & C1634 & B & \texttt{Device:\allowbreak{}C} & \textbf{DNP} CFF across R954 \\
C957, C958 & 22 µF & C\_\allowbreak{}1206 & 25 V X5R & C12891 & B & \texttt{Device:\allowbreak{}C} & VOUT \\
C959 & 100 nF & C\_\allowbreak{}0603 & 50 V & C14663 & B & \texttt{Device:\allowbreak{}C} & VOUT HF \\
D951 & Red LED KT-0603R & \texttt{LED\_\allowbreak{}SMD:\allowbreak{}LED\_\allowbreak{}0603\_\allowbreak{}1608Metric} & & C2286 & B & \texttt{Device:\allowbreak{}LED} & \textbf{DNP} power LED \\
R956 & 1.0 kΩ & R\_\allowbreak{}0603 & & C21190 & B & \texttt{Device:\allowbreak{}R} & \textbf{DNP} (with D951) \\
TP951 & +3V3 & \texttt{TestPoint:\allowbreak{}TestPoint\_\allowbreak{}Pad\_\allowbreak{}D1.5mm} & & -- & -- & \texttt{Connector:\allowbreak{}TestPoint} & \\
TP952, TP953, TP954 & GND & \texttt{TestPoint:\allowbreak{}TestPoint\_\allowbreak{}Pad\_\allowbreak{}D2.0mm} & & -- & -- & \texttt{Connector:\allowbreak{}TestPoint} & one near each converter \\
\end{longtable}
}

\subparagraph{4.2.2 U951 pin-by-pin}

{\def\LTcaptype{none} % do not increment counter
\begin{longtable}[]{@{}
  >{\raggedright\arraybackslash}p{(\linewidth - 6\tabcolsep) * \real{0.2500}}
  >{\raggedright\arraybackslash}p{(\linewidth - 6\tabcolsep) * \real{0.2500}}
  >{\raggedright\arraybackslash}p{(\linewidth - 6\tabcolsep) * \real{0.2500}}
  >{\raggedright\arraybackslash}p{(\linewidth - 6\tabcolsep) * \real{0.2500}}@{}}
\toprule\noalign{}
\begin{minipage}[b]{\linewidth}\raggedright
Pin
\end{minipage} & \begin{minipage}[b]{\linewidth}\raggedright
Name
\end{minipage} & \begin{minipage}[b]{\linewidth}\raggedright
Net
\end{minipage} & \begin{minipage}[b]{\linewidth}\raggedright
Connection
\end{minipage} \\
\midrule\noalign{}
\endhead
\bottomrule\noalign{}
\endlastfoot
1 & RT & no connect (KiCad no-connect flag) & Floating = 500 kHz (Table 9-1) \\
2 & EN & \texttt{BUCK\_\allowbreak{}EN} & R951 from \texttt{VSYS}, R952 to GND \\
3 & VIN & \texttt{VSYS} & C951--C953 \\
4 & GND & \texttt{GND} & \\
5 & SW & \texttt{BUCK\_\allowbreak{}SW} & L951 pin 1; C955 pin 2 \\
6 & BST & \texttt{BUCK\_\allowbreak{}BST} & R953 → \texttt{BUCK\_\allowbreak{}BST\_\allowbreak{}C} → C955 → \texttt{BUCK\_\allowbreak{}SW} \\
7 & SS & \texttt{BUCK\_\allowbreak{}SS} & C954 33 nF to GND \\
8 & FB & \texttt{BUCK\_\allowbreak{}FB} & R954 from \texttt{+3V3}, R955 to GND, C956 (DNP) from \texttt{+3V3} \\
\end{longtable}
}

\textbf{L951}: pin 1 \texttt{BUCK\_\allowbreak{}SW}, pin 2 \texttt{+3V3}. \textbf{D951/R956 (DNP)}: \texttt{+3V3} → R956 → D951 anode; D951 cathode → GND.

\subparagraph{4.2.3 Calculations (TPS62933 SLUSEA4D)}

\begin{itemize}
\tightlist
\item
  \textbf{VOUT} (eq. 14): 0.8 × (1 + 75/24) = \textbf{3.300 V} (VFB 0.792--0.808 V at 25 °C → 3.27--3.33 V).
\item
  \textbf{UVLO} (§9.3.6, eq. 3--5; VEN\_\allowbreak{}RISE 1.21 V, VEN\_\allowbreak{}FALL 1.17 V, Ip 0.7 µA, Ih 1.4 µA).

  \begin{itemize}
  \tightlist
  \item
    R1 = 300 k, R2 = 82 k give start = (1.21 − 0.7 µ × 64.4 k) × 382/82 = \textbf{5.43 V} and stop = (1.17 − 2.1 µ × 64.4 k) × 4.659 = \textbf{4.82 V}. Hysteresis 0.61 V ≥ 0.5 V as recommended.
  \item
    VEN = 2.07 V at 9 V in and 4.43 V even at a 20 V fault, both \textless{} 5.5 V max.
  \item
    The battery protector cuts at 6.0 V and VSYSMIN is 7.0 V with VBUS present, so 3V3 never chatters.
  \end{itemize}
\item
  \textbf{fsw} 500 kHz (RT floating). Ripple ΔI = 3.3 × (8.6 − 3.3)/(8.6 × 4.7 µ × 500 k) = 0.87 A. Peak at a 0.8 A load = 1.23 A \textless{} Isat 3 A.
\item
  \textbf{COUT}: 2 × 22 µF/25 V at 3.3 V ≈ 2 × 17 µF ≈ 34 µF effective. Table 10-2: typical 40 µF, minimum 15 µF.
\item
  \textbf{SS}: tss = 33 nF × 0.8 V/5.5 µA = \textbf{4.8 ms} (§10.2.2.4).
\item
  \textbf{Inrush and sequencing.} +3V3 comes up whenever VSYS \textgreater{} 5.43 V, so it is always on with VBUS (VSYS ≥ 7 V). Ship mode turns it off by removing VSYS.
\end{itemize}

\begin{center}\rule{0.5\linewidth}{0.5pt}\end{center}

\subsubsection{5. Test points (summary)}

{\def\LTcaptype{none} % do not increment counter
\begin{longtable}[]{@{}lll@{}}
\toprule\noalign{}
TP & Net & Sheet \\
\midrule\noalign{}
\endhead
\bottomrule\noalign{}
\endlastfoot
TP601 & \texttt{VBUS} & usb\_\allowbreak{}power \\
TP602 & \texttt{PD\_\allowbreak{}PG} (optional) & usb\_\allowbreak{}power \\
TP701 & \texttt{VSYS} & charger \\
TP702 & \texttt{CHG\_\allowbreak{}BAT} (optional) & charger \\
TP801 & \texttt{VBAT\_\allowbreak{}PACK} & battery \\
TP802 & \texttt{BATT\_\allowbreak{}N} & battery \\
TP803 & \texttt{BAL\_\allowbreak{}MID} & battery \\
TP901 & \texttt{PVDD} & rails \\
TP951 & \texttt{+3V3} & rails \\
TP952--TP954 & \texttt{GND} & rails \\
\end{longtable}
}

Indicators: D701 (charge STAT, fitted) and D951 (3V3 power, DNP to save standby current).

\begin{center}\rule{0.5\linewidth}{0.5pt}\end{center}

\subsubsection{6. Layout checklist}

\textbf{Charger (U701)}

\begin{itemize}
\tightlist
\item
  Follow the TI EVM (SLUUCB5E §4).
\item
  PMID caps (C707 first, then C704--C706) on the same layer, straight across PMID (29) ↔ GND (27). This is the buck-mode hot loop through Q1/Q2.
\item
  SYS caps (C716 first) straight across SYS (25) ↔ GND. This is the boost-mode hot loop through Q3/Q4.
\item
  Keep the SW1 (28) and SW2 (26) copper just large enough for L701.
\item
  BTST caps go within 1 mm of their pins.
\item
  REGN cap at pin 5.
\item
  GND pin 27 gets a via array into the inner GND plane.
\item
  PROG/ILIM/TS resistors return to a quiet GND spot, away from SW nodes.
\end{itemize}

\textbf{Boost (U901)}

\begin{itemize}
\tightlist
\item
  Hot loop SW (7) → VOUT (8) → C914 / first C906 → PGND (6) under 5 mm, on the top layer. The bars are 0.55 mm apart, so put C914 directly across pins 8--6.
\item
  Wide pours on VSYS/SW/PVDD, with thermal vias under the PGND bar and the GND pins. Up to ≈ 2 W of IC loss at 40 W.
\item
  AGND (2) connects to PGND only at the pin-6 pad.
\item
  FB divider (R904--R907) within 3 mm of pin 5. Route the R904 sense from the C906 PVDD node as a thin trace away from L901/SW.
\item
  COMP parts next to pin 4.
\end{itemize}

\textbf{Buck (U951)}

\begin{itemize}
\tightlist
\item
  C953/C951 across VIN (3)--GND (4) within 1 mm. SW node small. FB parts at pin 8, away from SW and L951.
\end{itemize}

\textbf{Protection and pack (high current, design for 7 A continuous / 10 A peak)}

\begin{itemize}
\tightlist
\item
  \texttt{PACK\_\allowbreak{}P} → F801 → \texttt{VBAT\_\allowbreak{}PACK} → Q701 → \texttt{CHG\_\allowbreak{}BAT}, and \texttt{BATT\_\allowbreak{}N} → Q801 → Q802 → \texttt{GND}, need ≥ 3 mm traces on 1 oz plus duplicates on inner/bottom layers, or pours.
\item
  \textbf{\texttt{BATT\_\allowbreak{}N} is a separate copper island.} It must never touch the GND plane or GND pours (DRC: keep a net-class clearance). Only J801 pin 1, J802 via JP801, Q801 source, U801 VSS, C801/C802/C803, U802/Q803, R811 and TP802 connect to it.
\item
  R803 (V−) returns at the Q802 source / GND junction.
\item
  The common-drain \texttt{PROT\_\allowbreak{}D} gets a pour joining both drain tabs.
\item
  \textbf{Q801 needs copper for 1.5 W} (body-diode conduction during over-discharge recovery at the 2 A default, S3). Use ≥ 2 cm² on the drain tab and the BATT\_\allowbreak{}N source, with thermal vias into the matching islands (never into GND for BATT\_\allowbreak{}N). Keep it away from the pack NTC lead and TH801.
\item
  Q701: leave room next to it for the DNP C720 (gate--drain) and C721 (gate--source) at SFET\_\allowbreak{}G, plus a test pad on SFET\_\allowbreak{}G (ECO-002 \#5).
\end{itemize}

\textbf{Balancing}

\begin{itemize}
\tightlist
\item
  R805/R806/R808/R809 dissipate up to 0.3 W per cell. Spread them out, keep them away from TH801 and U801, and give them copper for heat.
\end{itemize}

\textbf{NTC}

\begin{itemize}
\tightlist
\item
  J803 (the default pack NTC) goes at the pack-lead exit next to J801/J802 (S6). TH801 (alternate) sits next to J801, away from the boost, charger and amps. JP803 is reachable for re-jumpering. The \texttt{BAT\_\allowbreak{}TS} trace runs away from SW nodes, with a GND reference.
\end{itemize}

\textbf{USB-C}

\begin{itemize}
\tightlist
\item
  D601, U601 and U602 within 5 mm of J601.
\item
  D+/D− as a 90 Ω differential pair (short; full speed is tolerant).
\item
  U603 (0.4 mm DFN) EP gets ≥ 2 vias; C604 at pin 1.
\item
  Q601/Q602 next to U603.
\end{itemize}

\textbf{Thermal}

\begin{itemize}
\tightlist
\item
  Heat sources: U901 + L901 (≈ 0.5 W typical, ≈ 2.5 W worst), U701 during charge-while-play (≈ 1.5 W), the battery path (≈ 1 W at 6 A), Q801 during over-discharge recovery (≈ 1.5 W, S3), the balancers (≈ 0.6 W, end of charge only). Spread them across the board.
\end{itemize}

\textbf{Bulk}

\begin{itemize}
\tightlist
\item
  C920 (10 × 10.2 mm electrolytic) goes in mech zone D (≤ 24 mm) between the boost output and the amps' PVDD entry, on a wide PVDD pour (≥ 3 A). The amp bulks C301/C401 go in zone G or on the bottom side (S11).
\end{itemize}

\begin{center}\rule{0.5\linewidth}{0.5pt}\end{center}

\subsubsection{7. Firmware notes (for the firmware spec)}

\begin{enumerate}
\def\labelenumi{\arabic{enumi}.}
\tightlist
\item
  \textbf{BQ25792 (0x6B): protection contract (release review S1, ECO-002 \#7). Mandatory; do this first at every boot.}

  \begin{itemize}
  \tightlist
  \item
    Write \textbf{SFET\_\allowbreak{}PRESENT = 1, EN\_\allowbreak{}BATOC = 1, WATCHDOG = 0} (watchdog disabled).
  \item
    If the watchdog must stay enabled, re-arm all three and ICHG on every WD\_\allowbreak{}FLAG/INT, because \textbf{a watchdog expiry or REG\_\allowbreak{}RST resets EN\_\allowbreak{}BATOC to 0 and ICHG to 2 A} (BQ25798 Table 7-30; BQ25792 Table 9-2).
  \item
    \textbf{Read back} and verify. If it fails, keep BOOST\_\allowbreak{}EN low and report a fault.
  \item
    Bench acceptance: 20 mΩ PVDD short with blank firmware and with armed firmware (§4.1.4).
  \item
    Then:

    \begin{itemize}
    \tightlist
    \item
      Set AUTO\_\allowbreak{}INDET\_\allowbreak{}EN = 0 (D+/D− are not connected).
    \item
      \textbf{Write ICHG at boot}: ≤ 2.5 A for 5 Ah, ≤ 1.75 A for 3.5 Ah. The POR default is 2 A.
    \item
      VREG 8.4 V (or 8.3 V for longevity).
    \item
      IINDPM from CH224Q reg 0x50 × 50 mA, capped at 3.0 A. With no PD (CH224Q reg 0x09), use ≤ 0.5--1.5 A (L1).
    \item
      JEITA: stop charging above 45 °C (ECO-001 \#6).
    \item
      EN\_\allowbreak{}IBAT (ADC) for coulomb counting.
    \item
      Read TS in battery mode and stop playback above the pack limit (no hardware discharge cut-off, S6).
    \end{itemize}
  \item
    Ship mode = SDRV\_\allowbreak{}CTRL 10b (with no adapter). Wake with QON: SW701 or the UI button, ≥ 1 s. A ≥ 10 s press resets the system.
  \item
    I2C: run 100--200 kHz until the rise time is scoped with VBUS present and absent (S5).
  \end{itemize}
\item
  \textbf{CH224Q (0x22/0x23)}: only reachable while VBUS is present. Read reg 0x09 (protocol), 0x50 (current) and 0x60--0x8F (source PDOs). Optional: write 0x0A to change the request. Never select 28 V/EPR.
\item
  \textbf{Boost}

  \begin{itemize}
  \tightlist
  \item
    Power-up:

    \begin{enumerate}
    \def\labelenumii{\arabic{enumii}.}
    \tightlist
    \item
      \texttt{PVDD\_\allowbreak{}TRIM} = 3.3 V for ≥ 150 ms.
    \item
      BOOST\_\allowbreak{}EN high.
    \item
      Wait \textbf{≥ 20 ms}, then release amp PDN.
    \item
      Ramp \texttt{PVDD\_\allowbreak{}TRIM} to the target over ≥ 30--50 ms.
    \end{enumerate}

    This soft start avoids the 19.4--21 V start-up overshoot (§4.1.4).
  \item
    Power-down: PDN low first, then BOOST\_\allowbreak{}EN low.
  \item
    PVDD\_\allowbreak{}TRIM (R907/D901 fitted): 1.5--3.3 V → 18.0--12.2 V (SPICE curve in §4.1.4). Use LEDC ≥ 20 kHz or DAC1. Allow ≥ 150 ms after each trim change.
  \item
    Keep the amp DRC/limiter so the average battery current stays ≤ 6 A.
  \item
    With BOOST\_\allowbreak{}EN low, PVDD is ≈ VSYS − 0.5 V (not 0 V) because there is no disconnect FET. Keep amp PDN low whenever the boost is off. The TAS5805M register setup can happen with PVDD at either level.
  \item
    EN\_\allowbreak{}BATOC = 1 is also the PVDD short-circuit protection (§4.1.4).
  \end{itemize}
\end{enumerate}
